\documentclass[../main.tex]{subfiles}
\graphicspath{{\subfix{../image/}}{\subfix{../figure/}}}

\begin{document}

\chapter{有限 Radon 测度序列}\label{chapter:ch05}

设 $X$ 为局部紧 Hausdorff 空间。$X$ 上有限带号 Radon 测度所成的向量空间 $\mathscr{R}_f(X)$，赋以范数 $\|\mu\|:=|\mu|(X)$（$|\mu|$ 为测度 $\mu$ 的全变差），是一个 Banach 空间；事实上，它是 $X$ 上有限带号测度所成 Banach 空间 $(\mathcal{M}_f(X),\|\cdot\|)$ 的闭子空间。本章首先讨论该空间中不同类型的收敛，随后建立它们与 Young 测度相应收敛概念之间的关系。在此之后，我们对 $L^1(X)$ 中有界序列的收敛得出结论。

\section{收敛的概念}\label{section:ch05-Radon弱星}

\subsection{强收敛}

\cref{proposition:ch02-强收敛性质}启发我们给出以下定义，特别地，当 $X$ 为赋 Borel $\sigma$-代数的局部紧 Hausdorff 空间时，该定义适用于 $\mathscr{R}_f(X)$。

\begin{definition}\label{definition:ch05-强收敛}
设 $(X,\mathcal{A})$ 为可测空间。序列 $\{\mu_k\}\subseteq\mathcal{M}_f(X)$ \textbf{强收敛}于 $\mu\in\mathcal{M}_f(X)$（记作 $\mu_k\to\mu$），若
\begin{align*}
\lim_{k\to\infty}\|\mu_k-\mu\|
=\lim_{k\to\infty}|\mu_k-\mu|(X)=0.
\end{align*}
$\mathcal{M}_f(X)$ 上的强拓扑为与范数 $\|\cdot\|$ 相伴的度量拓扑。为清楚起见，我们将上述范数在子空间 $\mathscr{R}_f(X)\subseteq\mathcal{M}_f(X)$ 上的限制记作 $\|\cdot\|_{\mathscr{R}_f(X)}$。
\end{definition}

\subsection{弱$^*$收敛}

\cref{theorem:ch02-Rf对偶}启发我们给出以下定义。

\begin{definition}\label{definition:ch05-弱星收敛}
设 $X$ 为局部紧 Hausdorff 空间。序列 $\{\mu_k\}\subseteq\mathscr{R}_f(X)$ \textbf{弱$^*$收敛}于 $\mu\in\mathscr{R}_f(X)$（记作 $\mu_k\stackrel{*}{\rightharpoonup}\mu$），若
\begin{align}
\lim_{k\to\infty}\int_X g\,\mathrm{d}\mu_k=\int_X g\,\mathrm{d}\mu
\quad\text{对任意 }g\in C_0(X).\label{align:ch05-5.1}
\end{align}
$\mathscr{R}_f(X)$ 上的弱$^*$拓扑，记作 $\mathcal{T}_{w^*}$，是使所有映射 $\mu\mapsto\int_X g\,\mathrm{d}\mu$（$g\in C_0(X)$）连续的最弱拓扑。
\end{definition}

\begin{remark}\label{remark:ch05-弱星积分}
由于 $C_0(X)\subseteq C_b(X)\subseteq L^1(X,\mathscr{B}(X),|\mu|)$ 且对任意 $k\in\mathbb{N}$ 有 $C_0(X)\subseteq C_b(X)\subseteq L^1(X,\mathscr{B}(X),|\mu_k|)$，故\eqref{align:ch05-5.1}与下文\eqref{align:ch05-5.3}中的积分是良定义的（见\cref{definition:ch02-Lebesgue}）。
\end{remark}

$\mathcal{T}_{w^*}$ 的一组基为 $\mathscr{R}_f(X)$ 的如下子集族：
\begin{align}
B_{g_1,\dots,g_n;\epsilon}(\mu_0)
:=\left\{\mu\in\mathscr{R}_f(X)\ \middle|\
\left|\int_X g_j\,\mathrm{d}\mu-\int_X g_j\,\mathrm{d}\mu_0\right|<\epsilon
\ \forall j=1,\dots,n\right\},\label{align:ch05-5.2}
\end{align}
其中 $g_1,\dots,g_n\in C_0(X)$，$\mu_0\in\mathscr{R}_f(X)$，$n\in\mathbb{N}$，$\epsilon>0$。

在对 $X$ 的适当假设下，限制在 $\mathscr{R}_f(X)$ 有界子集上的弱$^*$拓扑 $\mathcal{T}_{w^*}$ 是可度量化的：

\begin{proposition}\label{proposition:ch05-弱星可度量化}
\label{theorem:ch05-弱星可测}
设 $X$ 为具可数基的局部紧 Hausdorff 空间，$M\subseteq\mathscr{R}_f(X)$ 有界。则相对弱$^*$拓扑 $\mathcal{T}_{w^*}\cap M$ 可度量化。
\end{proposition}

\begin{proof}
空间 $C_0(X)$ 赋以 $\|\cdot\|_\infty$ 范数是可分的（见附录 A,\cref{subsection:appA-可分})。于是由\cref{theorem:ch02-Rf对偶}与\cref{proposition:appA-弱星可度量化}即得结论。
\end{proof}

\subsection{窄收敛}

在\cref{definition:ch05-弱星收敛}中选取不同的试验函数类，便得到 $\mathscr{R}_f(X)$ 上另一种收敛概念。

\begin{definition}\label{definition:ch05-窄收敛}
设 $X$ 为局部紧 Hausdorff 空间。序列 $\{\mu_k\}\subseteq\mathscr{R}_f(X)$ \textbf{窄收敛}于 $\mu\in\mathscr{R}_f(X)$（记作 $\mu_k\stackrel{n}{\rightharpoonup}\mu$），若
\begin{align}
\lim_{k\to\infty}\int_X g\,\mathrm{d}\mu_k=\int_X g\,\mathrm{d}\mu
\quad\text{对任意 }g\in C_b(X).\label{align:ch05-5.3}
\end{align}
$\mathscr{R}_f(X)$ 上的窄拓扑，记作 $\mathcal{T}_n$，是使所有映射 $\mu\mapsto\int_X g\,\mathrm{d}\mu$（$g\in C_b(X)$）连续的最弱拓扑。
\end{definition}

窄拓扑被某些作者称为“弱拓扑”（例如见文献~\cite{ref45}）。$\mathcal{T}_n$ 的一组基为\eqref{align:ch05-5.2}所定义的子集族 $B_{g_1,\dots,g_n;\epsilon}(\mu_0)\subseteq\mathscr{R}_f(X)$，只是其中 $g_1,\dots,g_n\in C_b(X)$。

\begin{remark}\label{remark:ch05-模糊收敛}
序列 $\{\mu_k\}\subseteq\mathscr{R}_f(X)$ \textbf{模糊收敛}于 $\mu\in\mathscr{R}_f(X)$，定义为
\begin{align*}
\lim_{k\to\infty}\int_X g\,\mathrm{d}\mu_k=\int_X g\,\mathrm{d}\mu
\quad\text{对任意 }g\in C_c(X)
\end{align*}
（例如见文献~\cite{ref42}）。显然，在 $\mathscr{R}_f(X)$ 中，强收敛 $\Rightarrow$ 窄收敛 $\Rightarrow$ 弱$^*$收敛 $\Rightarrow$ 模糊收敛。
\end{remark}

\begin{remark}\label{remark:ch05-范数下半连续}
\label{remark:ch05-5.1.3}
由范数的下半连续性可知，若 $\mu_k\stackrel{*}{\rightharpoonup}\mu$（特别地，若 $\mu_k\stackrel{n}{\rightharpoonup}\mu$），则
\begin{align}
|\mu|(X)=\|\mu\|_{\mathscr{R}_f(X)}
\leqslant\liminf_{k\to\infty}\|\mu_k\|_{\mathscr{R}_f(X)}
=\liminf_{k\to\infty}|\mu_k|(X).\label{align:ch05-5.4}
\end{align}
特别地，若 $\{\mu_k\}\subseteq\mathscr{R}^+_f(X)$，则 $\mu(X)\leqslant\liminf_{k\to\infty}\mu_k(X)$。具体地，若 $\mu_k\stackrel{n}{\rightharpoonup}\mu$，则在\eqref{align:ch05-5.3}中取 $g=1$ 得 $\lim_{k\to\infty}\mu_k(X)=\mu(X)$，故“没有质量损失”。若 $\mu_k\stackrel{*}{\rightharpoonup}\mu$，则未必如此：取 $X=\mathbb{R}$，$\mu_k=\delta_k$（$k\in\mathbb{N}$），则 $\mu_k\stackrel{*}{\rightharpoonup}0$ 且对所有 $k$ 有 $\mu_k(\mathbb{R})=1$。故在此情形不等式\eqref{align:ch05-5.4}是严格的。类似现象在函数序列中已遇到过（见\cref{example:ch02-弱收敛振荡}(i)）。
\end{remark}

\begin{remark}\label{remark:ch05-函数转测度}
设 $X$ 为局部紧 Hausdorff 空间，$\mu\in\mathscr{R}^+_f(X)$。设 $\{h_k\}\subseteq L^1(X,\mathscr{B}(X),\mu)$，对任意 $k\in\mathbb{N}$ 置
\begin{align}
\nu_k(E):=\int_E h_k\,\mathrm{d}\mu\qquad(E\in\mathscr{B}(X)).\label{align:ch05-5.5}
\end{align}
显然 $\{\nu_k\}\subseteq\mathscr{R}_f(X)$，且对所有 $k\in\mathbb{N}$ 有 $\nu_k\ll\mu$。$\{h_k\}$ 在 $L^1(X,\mathscr{B}(X),\mu)$ 中的收敛与 $\{\nu_k\}$ 在 $\mathscr{R}_f(X)$ 中的收敛之间存在明显联系：
\begin{itemize}
\item 若 $h_k\to h$ 在 $L^1(X)$ 中强收敛，且 $\nu(E):=\int_E h\,\mathrm{d}\mu$（$E\in\mathscr{B}(X)$），则 $\|\nu_k-\nu\|_{\mathscr{R}_f(X)}=|\nu_k-\nu|(X)=\int_X|h_k-h|\,\mathrm{d}\mu\to0$，故 $\nu_k\to\nu$ 在 $\mathscr{R}_f(X)$ 中强收敛；
\item 若 $h_k\to h$ 在 $L^1(X)$ 中弱收敛，则特别地
\begin{align*}
\lim_{k\to\infty}\int_X g h_k\,\mathrm{d}\mu=\int_X g h\,\mathrm{d}\mu
\quad\text{对任意 }g\in C_b(X),
\end{align*}
从而
\begin{align*}
\lim_{k\to\infty}\int_X g\,\mathrm{d}\nu_k=\int_X g\,\mathrm{d}\nu
\quad\text{对任意 }g\in C_b(X)
\end{align*}
其中 $\nu$ 同上。故 $\nu_k\stackrel{n}{\rightharpoonup}\nu$。
\end{itemize}
更一般地，即使序列 $\{h_k\}$ 在 $L^1(X)$ 中不收敛，由\eqref{align:ch05-5.5}定义的序列 $\{\nu_k\}$ 仍可（弱$^*$或窄）收敛于某个 $\nu\in\mathscr{R}_f(X)$（见\cref{example:ch05-测度逼近}(i)）。在此类情形中，按记号滥用，我们偶尔写作 $h_k\stackrel{*}{\rightharpoonup}\nu$ 或 $h_k\stackrel{n}{\rightharpoonup}\nu$。
\end{remark}

牢记\cref{remark:ch05-函数转测度}，我们指出一些与 $\mathscr{R}_f(X)$ 中收敛相关的现象。

\begin{example}\label{example:ch05-测度逼近}
\label{example:ch05-振荡集中}
(i) 每个 $\nu\in\mathscr{R}^+_f(\mathbb{R})$ 都是序列 $\{h_k\}\subseteq C^\infty_c(\mathbb{R})$（在 $L^1(\mathbb{R})\equiv L^1(\mathbb{R},\mathscr{B},\lambda)$ 中有界）的弱$^*$极限（更确切地说，是\eqref{align:ch05-5.5}中定义的密度 $h_k$ 所对应的测度序列 $\nu_k$ 的弱$^*$极限）。事实上，设 $\rho_{1/k}$ 为磨光子，$\rho_{1/k}(x):=k\rho(kx)$，其中 $\rho\in C^\infty_c(\mathbb{R})$，$\rho\geqslant0$，$\int_\mathbb{R}\rho\,\mathrm{d}\lambda=1$（$k\in\mathbb{N}$）。对任意 $k\in\mathbb{N}$，考虑卷积 $\nu*\rho_{1/k}\in C^\infty_c(\mathbb{R})$，
\begin{align*}
\int_\mathbb{R}g(x)(\nu*\rho_{1/k})(x)\,\mathrm{d}\lambda(x)
:=\int_\mathbb{R}\mathrm{d}\nu(x)\int_\mathbb{R}\rho_{1/k}(y-x)g(y)\,\mathrm{d}\lambda(y)
\qquad(g\in C^\infty_c(\mathbb{R}))
\end{align*}
（见\cref{definition:ch03-磨光}(ii)与\cref{remark:ch03-磨光注}）。置 $h_k:=\nu*\rho_{1/k}$，$\nu_k(E):=\int_E h_k\,\mathrm{d}\lambda$（$E\in\mathscr{B}$）。由\cref{proposition:ch03-磨光性质}(i)，
\begin{align*}
\lim_{k\to\infty}\int_\mathbb{R}g\,\mathrm{d}\nu_k
=\lim_{k\to\infty}\int_\mathbb{R}g h_k\,\mathrm{d}\lambda
=\int_\mathbb{R}g\,\mathrm{d}\nu
\quad\text{对任意 }g\in C^\infty_c(\mathbb{R}).
\end{align*}
易见 $\|h_k\|_{L^1(\mathbb{R})}\leqslant\|\nu\|=\nu(X)<\infty$，故由标准论证上述等式对任意 $g\in C_0(\mathbb{R})$ 成立。

(ii) 特别地，上述考虑适用于 $\nu$ 为原子测度的情形，例如 $\nu=\delta_0$。因此序列 $\{\nu_k\}\subseteq\mathscr{R}^+_f(\mathbb{R})$（对所有 $k\in\mathbb{N}$，$\nu_k\ll\lambda$）可弱$^*$收敛于测度 $\nu\perp\lambda$。在此情形我们也说发生了质量集中（见\cref{example:ch02-弱收敛振荡}）。反之，考虑测度空间 $(I,\mathscr{B}\cap I,\lambda|_{\mathscr{B}\cap I})$，其中 $I\equiv(0,1)$。对任意固定的 $h\in C_b(I)$，置 $\nu(E):=\int_E h\,\mathrm{d}\lambda$（$E\in\mathscr{B}\cap I$），并考虑原子测度序列
\begin{align*}
\nu_k:=\frac{1}{k}\sum_{j=1}^k h\!\left(\frac{j}{k}\right)\delta_{j/k}\qquad(k\in\mathbb{N}).
\end{align*}
由初等结果，
\begin{align*}
\lim_{k\to\infty}\int_I g\,\mathrm{d}\nu_k
=\lim_{k\to\infty}\frac{1}{k}\sum_{j=1}^k g\!\left(\frac{j}{k}\right)h\!\left(\frac{j}{k}\right)
=\int_I gh\,\mathrm{d}\lambda=\int_I g\,\mathrm{d}\nu
\quad\text{对任意 }g\in C_b(I).
\end{align*}
故在此情形，对所有 $k\in\mathbb{N}$ 有 $\nu_k\perp\lambda$，$\nu\ll\lambda$，且 $\nu_k\stackrel{n}{\rightharpoonup}\nu$。
\end{example}

此后我们主要处理有限（正）Radon 测度序列。有限带号 Radon 测度序列的结果可由下文所证结果容易地得到（见\cref{remark:ch02-带号约化}）。

窄收敛可如下刻画。

\begin{theorem}[Portmanteau 定理]\label{theorem:ch05-portmanteau}
设 $X$ 为具可数基的局部紧 Hausdorff 空间，$\{\mu_k\}\subseteq\mathscr{R}^+_f(X)$，$\mu\in\mathscr{R}^+_f(X)$。则下列命题等价：
\begin{enumerate}[(i)]
\item\label{item:ch05-port-i} $\mu_k\stackrel{n}{\rightharpoonup}\mu$；
\item\label{item:ch05-port-ii} 对每个一致连续的 $g\in C_b(X)$，
\begin{align}
\lim_{k\to\infty}\int_X g\,\mathrm{d}\mu_k=\int_X g\,\mathrm{d}\mu;\label{align:ch05-5.6}
\end{align}
\item\label{item:ch05-port-iii} $\mu_k(X)\to\mu(X)$，且对任意闭集 $C\subseteq X$，
\begin{align}
\limsup_{k\to\infty}\mu_k(C)\leqslant\mu(C);\label{align:ch05-5.7}
\end{align}
\item\label{item:ch05-port-iv} $\mu_k(X)\to\mu(X)$，且对任意开集 $A\subseteq X$，
\begin{align}
\liminf_{k\to\infty}\mu_k(A)\geqslant\mu(A);\label{align:ch05-5.8}
\end{align}
\item\label{item:ch05-port-v} 对所有满足 $\mu(\partial E)=0$ 的 $E\in\mathscr{B}(X)$，
\begin{align}
\lim_{k\to\infty}\mu_k(E)=\mu(E).\label{align:ch05-5.9}
\end{align}
\end{enumerate}
\end{theorem}

\begin{proof}
(i)$\Rightarrow$(ii)。由\cref{definition:ch05-窄收敛}显然。

(ii)$\Rightarrow$(iii)。取 $g=1$ 表明 $\mu_k(X)\to\mu(X)$。设 $X$ 赋以与给定拓扑相容的度量 $d$。若 $C=\emptyset$，不等式\eqref{align:ch05-5.7}显然成立。否则，对任意 $n\in\mathbb{N}$，置 $C_{1/n}:=\{x\in X\mid d(x,C)<1/n\}$。显然 $C_{1/n}$ 为开集，$C=\bigcap_{n=1}^\infty C_{1/n}$，故 $\mu(C)=\lim_{n\to\infty}\mu(C_{1/n})$。置 $f_n:X\to[0,\infty)$，$f_n(x):=\max\{1-nd(x,C),0\}$（$n\in\mathbb{N}$）。显然，对任意 $n\in\mathbb{N}$，函数 $f_n$ 在 $X$ 上一致连续，且 $\chi_C\leqslant f_n\leqslant\chi_{C_{1/n}}$。于是由\eqref{align:ch05-5.6}，对任意固定的 $n\in\mathbb{N}$，
\begin{align*}
\limsup_{k\to\infty}\mu_k(C)
\leqslant\lim_{k\to\infty}\int_X f_n\,\mathrm{d}\mu_k
=\int_X f_n\,\mathrm{d}\mu\leqslant\mu(C_{1/n}).
\end{align*}
令 $n\to\infty$ 即得\eqref{align:ch05-5.7}，结论得证。

(iii)$\Leftrightarrow$(iv)。由\eqref{align:ch05-5.7}，对任意开集 $A\subseteq X$，
\begin{align*}
\liminf_{k\to\infty}\mu_k(A)
=\lim_{k\to\infty}\mu_k(X)-\limsup_{k\to\infty}\mu_k(A^c)
\geqslant\mu(X)-\mu(A^c)=\mu(A),
\end{align*}
故\eqref{align:ch05-5.8}成立。反转论证即得逆命题。

(iv)$\Rightarrow$(v)。由假设 $\mu(\partial E)=\mu(\overline{E})-\mu(E^\circ)=0$，故 $\mu(E)=\mu(E^\circ)=\mu(\overline{E})$。于是由\eqref{align:ch05-5.7}–\eqref{align:ch05-5.8}，
\begin{align*}
\mu(E)=\mu(E^\circ)
\leqslant\liminf_{k\to\infty}\mu_k(E^\circ)
\leqslant\liminf_{k\to\infty}\mu_k(E)
\leqslant\limsup_{k\to\infty}\mu_k(E)
\leqslant\limsup_{k\to\infty}\mu_k(\overline{E})
\leqslant\mu(\overline{E})=\mu(E).
\end{align*}
故\eqref{align:ch05-5.9}成立。

(v)$\Rightarrow$(i)。对任意 $g\in C_b(X)$ 与 $t>0$，置 $E_t:=\{g>t\}$。由于 $g$ 连续，对 $\lambda$-几乎所有 $t\in(0,\infty)$，有 $\mu(\partial E_t)\leqslant\mu(\{g=t\})=0$，故由\eqref{align:ch05-5.9}，
\begin{align}
\lim_{k\to\infty}\mu_k(\{g>t\})=\mu(\{g>t\})
\quad\text{对 }\lambda\text{-几乎所有 }t\in(0,\infty).\label{align:ch05-5.10}
\end{align}
于是由等式\eqref{align:ch02-2.19}（取 $\alpha=1$）、\eqref{align:ch05-5.10}及控制收敛定理，对任意 $g\in C_b(X)$ 有
\begin{align*}
\lim_{k\to\infty}\int_X g\,\mathrm{d}\mu_k
=\lim_{k\to\infty}\int_{[0,\infty)}\mu_k(\{g>t\})\,\mathrm{d}\lambda(t)
=\int_{[0,\infty)}\mu(\{g>t\})\,\mathrm{d}\lambda(t)
=\int_X g\,\mathrm{d}\mu.
\end{align*}
证明完成。
\end{proof}

\begin{remark}\label{remark:ch05-窄收敛范数}
如\cref{remark:ch05-范数下半连续}已述，\cref{theorem:ch05-portmanteau}表明：若 $\{\mu_k\}\subseteq\mathscr{R}^+_f(X)$，$\mu\in\mathscr{R}^+_f(X)$，且 $\mu_k\stackrel{n}{\rightharpoonup}\mu$，则 $\|\mu_k\|_{\mathscr{R}_f(X)}\to\|\mu\|_{\mathscr{R}_f(X)}$。
\end{remark}

\subsection{Prokhorov 距离}

本小节结果对度量空间 $(X,d)$ 表述。然而由\cref{proposition:appA-可度量化}，当 $X$ 为具可数基的局部紧 Hausdorff 空间且 $d$ 为与其拓扑相容的度量时，这些结果同样成立。

设 $(X,d)$ 为度量空间。对任意 $E\subseteq X$ 与 $r>0$，置 $E_r:=\{x\in X\mid d(x,E)<r\}$。对任意 $\mu_1,\mu_2\in\mathscr{R}^+_f(X)$，定义
\begin{align}
d_P(\mu_1,\mu_2):=\inf\left\{r>0\ \middle|\
\mu_1(E)\leqslant\mu_2(E_r)+r,\ \mu_2(E)\leqslant\mu_1(E_r)+r
\ \forall E\in\mathscr{B}(X)\right\}.\label{align:ch05-5.11}
\end{align}
注意每个 $r\geqslant\max\{\|\mu_1\|_{\mathscr{R}_f(X)},\|\mu_2\|_{\mathscr{R}_f(X)}\}$ 都属于右端集合，故 $d_P(\mu_1,\mu_2)<\infty$。

\begin{proposition}\label{proposition:ch05-Prokhorov度量}
$(\mathscr{R}^+_f(X),d_P)$ 为度量空间。
\end{proposition}

\begin{proof}
我们证明映射 $(\mu_1,\mu_2)\mapsto d_P(\mu_1,\mu_2)$ 是 $\mathscr{R}^+_f(X)$ 上的度量。显然 $d_P(\mu_1,\mu_2)=d_P(\mu_2,\mu_1)$。此外：
\begin{enumerate}[(i)]
\item 对所有 $E\in\mathscr{B}(X)$，有 $E\subseteq E_r$，故对所有 $\mu\in\mathscr{R}^+_f(X)$ 与 $r>0$ 有 $\mu(E)\leqslant\mu(E_r)$。因此
\begin{align*}
d_P(\mu,\mu):=\inf\{r>0\mid\mu(E)\leqslant\mu(E_r)+r\ \forall E\in\mathscr{B}(X)\}=0;
\end{align*}
\item 若 $d_P(\mu_1,\mu_2)=0$，则存在序列 $r_n\to0^+$ 使对所有 $n$ 与 $E\in\mathscr{B}(X)$ 有 $\mu_1(E)\leqslant\mu_2(E_{r_n})+r_n$，$\mu_2(E)\leqslant\mu_1(E_{r_n})+r_n$。由于 $\bigcap_{n=1}^\infty\overline{E_{r_n}}=\overline{E}$，故对所有 $E\in\mathscr{B}(X)$ 有 $\mu_1(E)\leqslant\mu_2(E)$ 与 $\mu_2(E)\leqslant\mu_1(E)$，从而对所有闭集 $C\in\mathscr{B}(X)$ 有 $\mu_1(C)=\mu_2(C)$。特别地 $\mu_1(X)=\mu_2(X)$，故对任意开集 $A\in\mathscr{B}(X)$ 有 $\mu_1(A)=\mu_2(A)$。于是由 $\mu_1,\mu_2$ 的正则性（见\cref{lemma:ch01-正则性}(i)）得 $\mu_1=\mu_2$；
\item 设 $\mu_1,\mu_2,\mu_3\in\mathscr{R}^+_f(X)$，$r,s>0$ 使 $d_P(\mu_1,\mu_2)<r$，$d_P(\mu_2,\mu_3)<s$。则对所有 $E\in\mathscr{B}(X)$，
\begin{align*}
\mu_1(E)\leqslant\mu_2(E_r)+r\leqslant\mu_3(E_{r+s})+r+s,\qquad
\mu_3(E)\leqslant\mu_2(E_s)+s\leqslant\mu_1(E_{r+s})+r+s,
\end{align*}
故 $d_P(\mu_1,\mu_3)\leqslant r+s$。对 $r,s$ 取inf即得
\begin{align*}
d_P(\mu_1,\mu_3)\leqslant d_P(\mu_1,\mu_2)+d_P(\mu_2,\mu_3),
\end{align*}
即三角不等式。结论得证。
\end{enumerate}
\end{proof}

\begin{definition}\label{definition:ch05-Prokhorov}
\eqref{align:ch05-5.11}中定义的 $\mathscr{R}^+_f(X)$ 上的度量 $d_P$ 称为 \textbf{Prokhorov 度量}。
\end{definition}

我们现在证明如下结果。

\begin{lemma}\label{lemma:ch05-Prokhorov对称}
设 $(X,d)$ 为度量空间，$\mu_1,\mu_2\in\mathscr{R}^+_f(X)$ 满足对所有 $E\in\mathscr{B}(X)$，
\begin{align}
\mu_1(E)\leqslant\mu_2(E_r)+r\qquad(r>0).\label{align:ch05-5.12}
\end{align}
则对所有 $F\in\mathscr{B}(X)$，
\begin{align}
\mu_2(F)\leqslant\mu_1(F_r)+r+\|\mu_2\|_{\mathscr{R}_f(X)}-\|\mu_1\|_{\mathscr{R}_f(X)}.\label{align:ch05-5.13}
\end{align}
\end{lemma}

\begin{remark}\label{remark:ch05-Prokhorov等范数}
由\cref{lemma:ch05-Prokhorov对称}，若\eqref{align:ch05-5.12}成立且 $\|\mu_1\|_{\mathscr{R}_f(X)}=\|\mu_2\|_{\mathscr{R}_f(X)}$，则
\begin{align}
d_P(\mu_1,\mu_2)
=\inf\{r>0\mid\mu_1(E)\leqslant\mu_2(E_r)+r\ \forall E\in\mathscr{B}(X)\}
=\inf\{r>0\mid\mu_2(E)\leqslant\mu_1(E_r)+r\ \forall E\in\mathscr{B}(X)\}.\label{align:ch05-5.14}
\end{align}
\end{remark}

\begin{proof}[证明\cref{lemma:ch05-Prokhorov}]
在\eqref{align:ch05-5.12}中取 $E=(F_r)^c$，立得
\begin{align*}
\mu_1(F_r)=\|\mu_1\|_{\mathscr{R}_f(X)}-\mu_1((F_r)^c)
=\|\mu_1\|_{\mathscr{R}_f(X)}-\mu_1(E)
\geqslant\|\mu_1\|_{\mathscr{R}_f(X)}-\mu_2(E_r)-r\\
=\|\mu_1\|_{\mathscr{R}_f(X)}-\|\mu_2\|_{\mathscr{R}_f(X)}
+\mu_2((E_r)^c)-r
\geqslant\|\mu_1\|_{\mathscr{R}_f(X)}-\|\mu_2\|_{\mathscr{R}_f(X)}
+\mu_2(F)-r.
\end{align*}
（注意 $E\subseteq(F_r)^c\Leftrightarrow F\subseteq(E_r)^c$。）
\end{proof}

\begin{example}\label{example:ch05-Delta距离}
设 $(X,d)$ 为度量空间。则对所有 $x_1,x_2\in X$，
\begin{align}
d_P(\delta_{x_1},\delta_{x_2})=\min\{1,d(x_1,x_2)\}.\label{align:ch05-5.15}
\end{align}
为证\eqref{align:ch05-5.15}，由\eqref{align:ch05-5.14}注意到
\begin{align*}
d_P(\delta_{x_1},\delta_{x_2})
=\inf\{r>0\mid\chi_E(x_1)\leqslant\chi_{E_r}(x_2)+r\ \forall E\in\mathscr{B}(X)\}.
\end{align*}
由于 $\chi_E(x_1)\leqslant1$，有 $d_P(\delta_{x_1},\delta_{x_2})\leqslant1$。此外，若 $r>d(x_1,x_2)$，则对所有 $E\in\mathscr{B}(X)$ 有 $\chi_E(x_1)\leqslant\chi_{E_r}(x_2)+r$：事实上，若 $x_1\in E$（否则不等式平凡）且 $r>d(x_1,x_2)$，则 $x_2\in E_r$，故 $\chi_{E_r}(x_2)=1$，不等式成立。综上 $d_P(\delta_{x_1},\delta_{x_2})\leqslant\min\{1,d(x_1,x_2)\}$。

为证反向不等式，取 $E=\{x_1\}$，故 $\chi_E(x_1)=1$。若 $d(x_1,x_2)\geqslant1$ 且 $r\in(0,1)$，则 $x_2\in(E_r)^c$，故 $\chi_{E_r}(x_2)=0$。于是
\begin{align}
\chi_E(x_1)=1>r=\chi_{E_r}(x_2)+r,\label{align:ch05-5.16}
\end{align}
故 $d_P(\delta_{x_1},\delta_{x_2})\geqslant1=\min\{1,d(x_1,x_2)\}$。类似地，若 $d(x_1,x_2)<1$ 且 $r\in(0,d(x_1,x_2)]$，则 $x_2\in(E_r)^c$，不等式\eqref{align:ch05-5.16}再次成立。故 $d_P(\delta_{x_1},\delta_{x_2})\geqslant d(x_1,x_2)=\min\{1,d(x_1,x_2)\}$。于是等式\eqref{align:ch05-5.15}成立。
\end{example}

\begin{remark}\label{remark:ch05-Delta等距}
由\cref{example:ch05-Delta距离}，映射 $X\to\mathscr{R}^+_f(X)$，$x\mapsto\delta_x$，是 $(X,\min\{1,d\})$ 到 $(\mathscr{R}^+_f(X),d_P)$ 的等距嵌入。注意 $\min\{1,d\}$ 是 $X$ 上与 $d$ 等价的度量。
\end{remark}

现在我们可以证明 $\mathscr{R}^+_f(X)$ 上的窄拓扑是可度量化的。

\begin{proposition}\label{proposition:ch05-窄拓扑可度量化}
设 $(X,d)$ 为可分度量空间。则 $\mathscr{R}^+_f(X)$ 上的窄拓扑与 Prokhorov 度量相伴的度量拓扑一致。
\end{proposition}

\begin{remark}\label{remark:ch05-Prokhorov推论}
由\cref{proposition:ch05-窄拓扑可度量化}，若 $(X,d)$ 为可分度量空间，则：
\begin{enumerate}[(i)]
\item\label{item:ch05-prok-i} $\mu_k\stackrel{n}{\rightharpoonup}\mu$ 当且仅当 $d_P(\mu_k,\mu)\to0$；
\item\label{item:ch05-prok-ii} 子集 $M\subseteq\mathscr{R}^+_f(X)$ 在窄拓扑中相对序列紧，当且仅当它在 $(\mathscr{R}^+_f(X),d_P)$ 中相对紧。
\end{enumerate}
\end{remark}

为证\cref{proposition:ch05-窄拓扑可度量化}，我们先建立如下引理。

\begin{lemma}\label{lemma:ch05-Prokhorov到窄}
设 $(X,d)$ 为度量空间，$\{\mu_k\}\subseteq\mathscr{R}^+_f(X)$，$\mu\in\mathscr{R}^+_f(X)$ 满足 $d_P(\mu_k,\mu)\to0$。则 $\mu_k\stackrel{n}{\rightharpoonup}\mu$。
\end{lemma}

\begin{proof}
设 $\{\epsilon_k\}\subseteq(0,\infty)$ 满足当 $k\to\infty$ 时 $\epsilon_k\to0^+$，且对所有 $k\in\mathbb{N}$ 有 $d_P(\mu_k,\mu)<\epsilon_k$。对任意闭集 $C\subseteq X$，置 $C_k\equiv C_{\epsilon_k}:=\{x\in X\mid d(x,C)<\epsilon_k\}$，故 $C=\bigcap_{k=1}^\infty C_k$ 且 $\mu(C)=\lim_{k\to\infty}\mu(C_k)$。由于 $d_P(\mu_k,\mu)<\epsilon_k$，由定义\eqref{align:ch05-5.11}，对任意 $k\in\mathbb{N}$ 有
\begin{align}
\mu_k(C)\leqslant\mu(C_k)+\epsilon_k,\qquad
\mu(C)\leqslant\mu_k(C_k)+\epsilon_k.\label{align:ch05-5.17}
\end{align}
由于 $\lim_{k\to\infty}\mu(C_k)=\mu(C)$ 且 $\lim_{k\to\infty}\epsilon_k=0$，在\eqref{align:ch05-5.17}第一式中令 $k\to\infty$ 得\eqref{align:ch05-5.7}。另一方面，在\eqref{align:ch05-5.17}第二式与\eqref{align:ch05-5.7}中均取 $C=X$，得
\begin{align*}
\mu(X)\leqslant\liminf_{k\to\infty}\mu_k(X)\leqslant\limsup_{k\to\infty}\mu_k(X)\leqslant\mu(X),
\end{align*}
故 $\lim_{k\to\infty}\mu_k(X)=\mu(X)$。于是由\cref{theorem:ch05-portmanteau}（见断言 (i) 与 (iii)）结论得证。
\end{proof}

\begin{proof}[证明\cref{proposition:ch05-窄拓扑可度量化}]
由\cref{lemma:ch05-Prokhorov到窄}，我们只需证明 $\mu_k\stackrel{n}{\rightharpoonup}\mu$ 蕴含 $d_P(\mu_k,\mu)\to0$。为此，设 $D\equiv\{x_k\}$ 为 $X$ 的可数稠密子集。

对任意固定的 $\epsilon>0$，置 $E_1:=B_{\epsilon/2}(x_1)$，$E_{m+1}:=B_{\epsilon/2}(x_{m+1})\setminus\bigcup_{k=1}^m E_k$（$m\in\mathbb{N}$）。显然，对所有 $m\in\mathbb{N}$ 有 $\operatorname{diam}(E_m)\leqslant\epsilon$，对所有 $l,m\in\mathbb{N}$，$l\neq m$ 有 $E_l\cap E_m=\emptyset$，且 $\bigcup_{m=1}^\infty E_m=X$。由于 $\mu(X)=\sum_{m=1}^\infty\mu(E_m)<\infty$，存在 $\bar{m}\in\mathbb{N}$ 使
\begin{align}
\mu\!\left(\bigcup_{m=\bar{m}+1}^\infty E_m\right)
=\sum_{m=\bar{m}+1}^\infty\mu(E_m)<\epsilon.\label{align:ch05-5.18}
\end{align}
考虑有限族
\begin{align*}
S:=\left\{U=E_{m_1}\cup E_{m_2}\cup\cdots\cup E_{m_p}\ \middle|\
1\leqslant m_1<m_2<\cdots<m_p\leqslant\bar{m}\right\}\subseteq\mathcal{P}(X),
\end{align*}
对任意 $U\in S$，置 $U_\epsilon:=\{x\in X\mid d(x,U)<\epsilon\}$。由于 $U_\epsilon$ 为开集且由假设 $\mu_k\stackrel{n}{\rightharpoonup}\mu$，由 Portmanteau 定理 $\liminf_{k\to\infty}\mu_k(U_\epsilon)\geqslant\mu(U_\epsilon)$。因此，由于 $S$ 有限，存在 $\bar{k}\in\mathbb{N}$ 使
\begin{align}
\mu_k(U_\epsilon)\geqslant\mu(U_\epsilon)-\epsilon
\quad\text{对所有 }k\geqslant\bar{k}\text{ 与 }U\in S.\label{align:ch05-5.19}
\end{align}
现注意到对所有 $F\in\mathscr{B}(X)$，
\begin{align}
F=\left[F\cap\bigcup_{m=1}^{\bar{m}}E_m\right]
\cup\left[F\cap\bigcup_{m=\bar{m}+1}^\infty E_m\right]
\subseteq[F\cap U_\epsilon]\cup\left[F\cap\bigcup_{m=\bar{m}+1}^\infty E_m\right]\label{align:ch05-5.20}
\end{align}
对某个 $U=E_{m_1}\cup E_{m_2}\cup\cdots\cup E_{m_p}\in S$（使对所有 $j=m_1,\dots,m_p$ 有 $E_{m_j}\cap F\neq\emptyset$）成立（注意这蕴含 $U_\epsilon\subseteq F_{2\epsilon}$，因为对任意 $m\in\mathbb{N}$ 有 $\operatorname{diam}(E_m)\leqslant\epsilon$）。由\eqref{align:ch05-5.18}–\eqref{align:ch05-5.20}得
\begin{align}
\mu(F)\leqslant\mu(U_\epsilon)+\epsilon
\leqslant\mu_k(U_\epsilon)+2\epsilon
\leqslant\mu_k(F_{2\epsilon})+2\epsilon
\quad\text{对所有 }k\geqslant\bar{k}.\label{align:ch05-5.21}
\end{align}
由 $F\in\mathscr{B}(X)$ 的任意性与\cref{lemma:ch05-Prokhorov对称}，从\eqref{align:ch05-5.21}得
\begin{align}
\mu_k(G)\leqslant\mu(G_{2\epsilon})+2\epsilon
+\|\mu_k\|_{\mathscr{R}_f(X)}-\|\mu\|_{\mathscr{R}_f(X)}
\leqslant\mu(G_{3\epsilon})+3\epsilon\label{align:ch05-5.22}
\end{align}
对所有 $G\in\mathscr{B}(X)$ 与所有充分大的 $k\in\mathbb{N}$ 成立，因为 $\lim_{k\to\infty}\|\mu_k\|_{\mathscr{R}_f(X)}=\|\mu\|_{\mathscr{R}_f(X)}$（见\cref{remark:ch05-窄收敛范数}）。

另一方面，由\eqref{align:ch05-5.21}，对所有充分大的 $k\in\mathbb{N}$ 与任意 $G\in\mathscr{B}(X)$，我们同样有
\begin{align}
\mu(G)\leqslant\mu_k(G_{3\epsilon})+3\epsilon.\label{align:ch05-5.23}
\end{align}
由定义\eqref{align:ch05-5.11}，从\eqref{align:ch05-5.22}–\eqref{align:ch05-5.23}可得：对所有 $\epsilon>0$ 与充分大的 $k\in\mathbb{N}$，$d_P(\mu_k,\mu)<\epsilon$。结论得证。
\end{proof}

\subsection{窄收敛与胎紧性}

刻画 $\mathscr{R}_f(X)$ 在弱$^*$与窄两种拓扑下的相对序列紧子集是有趣的。关于弱$^*$收敛，我们有如下结果：

\begin{theorem}\label{theorem:ch05-弱星紧性}
\label{theorem:ch05-弱星紧}
设 $X$ 为具可数基的局部紧 Hausdorff 空间，$M\subseteq\mathscr{R}_f(X)$。则下列命题等价：
\begin{enumerate}[(i)]
\item\label{item:ch05-wstar-i} $M$ 在弱$^*$拓扑中相对序列紧；
\item\label{item:ch05-wstar-ii} $M$ 有界。
\end{enumerate}
\end{theorem}

\begin{proof}
(i)$\Rightarrow$(ii)。由\cref{theorem:ch02-Rf对偶}，Banach 空间 $\mathscr{R}_f(X)$ 是 $C_0(X)$ 的对偶空间。若 $M$ 无界，则存在 $\{\mu_k\}\subseteq M$ 使 $\|\mu_k\|_{\mathscr{R}_f(X)}\to\infty$。另一方面，由假设，每个 $\{\mu_k\}\subseteq M$ 都含一个在 $\mathscr{R}_f(X)$ 中弱$^*$收敛的子列，因而由一致有界原理有界。矛盾证明了断言。

(ii)$\Rightarrow$(i) 由 Banach–Alaoglu 定理得到。
\end{proof}

关于窄收敛，我们从如下定义开始：

\begin{definition}\label{definition:ch05-胎紧}
设 $X$ 为局部紧 Hausdorff 空间。子集 $M\subseteq\mathscr{R}_f(X)$ 称为\textbf{胎紧的}（tight），若对任意 $\epsilon>0$，存在紧子集 $K\subseteq X$ 使对所有 $\mu\in M$ 有 $|\mu|(K^c)<\epsilon$。
\end{definition}

\begin{example}\label{example:ch05-胎紧}
设 $X=\mathbb{R}$，$\mu=\delta_{\bar{x}}$（$\bar{x}\in\mathbb{R}$）。故对所有 $E\in\mathscr{B}(\mathbb{R})$ 有 $\mu(E)=\chi_E(\bar{x})$。显然，对任意 $F\subseteq\mathbb{R}$，集合 $M=\{\delta_{\bar{x}}\mid\bar{x}\in F\}$ 胎紧当且仅当 $F$ 有界。注意若 $F$ 有界，则集合 $\{\alpha\delta_{\bar{x}}\mid\bar{x}\in F,\ \alpha>0\}$ 胎紧但无界。
\end{example}

为后续参考，我们证明如下结果。

\begin{proposition}\label{proposition:ch05-胎紧拓扑一致}
设 $X$ 为具可数基的局部紧 Hausdorff 空间，$M\subseteq\mathscr{R}^+_f(X)$ 胎紧。则 $M$ 上的窄拓扑与 $M$ 上的弱$^*$拓扑一致，即 $\mathcal{T}_n\cap M=\mathcal{T}_{w^*}\cap M$。
\end{proposition}

\begin{proof}
只需证明：对任意序列 $\{\mu_k\}\subseteq M$，若 $\mu_k\stackrel{*}{\rightharpoonup}\mu$，则 $\mu_k\stackrel{n}{\rightharpoonup}\mu$。固定 $\epsilon>0$，由\cref{definition:ch05-胎紧}存在紧子集 $K\subseteq X$ 使对所有 $k\in\mathbb{N}$ 有 $\mu_k(K^c)<\epsilon$。另一方面，由\cref{proposition:appA-Urysohn}，对任意开集 $A\subseteq X$，$A\supseteq K$，存在 $f\in C_c(X)$ 使 $f(X)\subseteq[0,1]$，$f|_K=1$，且 $\operatorname{supp}f\subseteq A$。固定 $g\in C_b(X)$。则 $fg\in C_0(X)$，且
\begin{align*}
\left|\int_X g\,\mathrm{d}\mu_k-\int_X g\,\mathrm{d}\mu\right|
\leqslant\left|\int_K g\,\mathrm{d}\mu_k-\int_K g\,\mathrm{d}\mu\right|
+\left|\int_{K^c}g\,\mathrm{d}\mu_k-\int_{K^c}g\,\mathrm{d}\mu\right|\\
\leqslant\left|\int_X fg\,\mathrm{d}\mu_k-\int_X fg\,\mathrm{d}\mu\right|
+\left|\int_X[\chi_K-f]g\,\mathrm{d}\mu_k\right|
+\left|\int_X[\chi_K-f]g\,\mathrm{d}\mu\right|
+\left|\int_{K^c}g\,\mathrm{d}\mu_k\right|
+\left|\int_{K^c}g\,\mathrm{d}\mu\right|
=:I_1+I_2+I_3+I_4+I_5.
\end{align*}
由于由假设 $\mu_k\stackrel{*}{\rightharpoonup}\mu$，存在 $\bar{k}\in\mathbb{N}$ 使对所有 $k>\bar{k}$，
\begin{align*}
I_1=\left|\int_X fg\,\mathrm{d}\mu_k-\int_X fg\,\mathrm{d}\mu\right|<\epsilon.
\end{align*}
我们还有，对所有 $k\in\mathbb{N}$，
\begin{align*}
I_2+I_4
=\left|\int_X[\chi_K-f]g\,\mathrm{d}\mu_k\right|
+\left|\int_{K^c}g\,\mathrm{d}\mu_k\right|
\leqslant 2\|g\|_\infty\mu_k(K^c)<2\|g\|_\infty\epsilon.
\end{align*}
此外，
\begin{align*}
I_3+I_5
=\left|\int_X[\chi_K-f]g\,\mathrm{d}\mu\right|
+\left|\int_{K^c}g\,\mathrm{d}\mu\right|
\leqslant 2\|g\|_\infty\mu(K^c)\leqslant 2\|g\|_\infty\epsilon,
\end{align*}
因为由范数的下半连续性（见\cref{remark:ch05-范数下半连续}），$\mu(K^c)\leqslant\liminf_{k\to\infty}\mu_k(K^c)$。

综上，我们证明了：对任意固定的 $g\in C_b(X)$ 与任意 $\epsilon>0$，存在 $\bar{k}\in\mathbb{N}$ 使对所有 $k>\bar{k}$，
\begin{align*}
\left|\int_X g\,\mathrm{d}\mu_k-\int_X g\,\mathrm{d}\mu\right|<(1+4\|g\|_\infty)\epsilon.
\end{align*}
结论得证。
\end{proof}

如下结果是\cref{theorem:ch05-弱星紧性}在弱$^*$收敛方面的对应物，表明窄收敛与胎紧性紧密交织。

\begin{theorem}[Prokhorov]\label{theorem:ch05-prokhorov}\label{lemma:ch05-Prokhorov}
设 $X$ 为具可数基的局部紧 Hausdorff 空间，$M\subseteq\mathscr{R}^+_f(X)$。则下列命题等价：
\begin{enumerate}[(i)]
\item\label{item:ch05-prokth-i} $M$ 在窄拓扑中相对序列紧；
\item\label{item:ch05-prokth-ii} $M$ 有界且胎紧。
\end{enumerate}
\end{theorem}

\begin{remark}\label{remark:ch05-polish}
更一般地，若 $X$ 为 Polish 空间（见附录 A,\cref{subsection:appA-Polish}），\cref{theorem:ch05-prokhorov}仍成立。\cref{proposition:ch05-ulam}亦然。
\end{remark}

\cref{theorem:ch05-prokhorov}的证明分两步给出，对应于\cref{proposition:ch05-ulam}与\cref{proposition:ch05-有界胎紧紧}。

\begin{proposition}[Ulam]\label{proposition:ch05-ulam}
设 $X$ 为具可数基的局部紧 Hausdorff 空间，$M\subseteq\mathscr{R}^+_f(X)$ 在窄拓扑中相对序列紧。则 $M$ 有界且胎紧。
\end{proposition}

\begin{proof}
由假设，对任意序列 $\{\mu_k\}\subseteq M$，存在 $\{\mu_{k_l}\}\subseteq\{\mu_k\}$ 与 $\mu\in\mathscr{R}_f(X)$ 使 $\mu_{k_l}\stackrel{n}{\rightharpoonup}\mu$，因而 $\|\mu_{k_l}\|_{\mathscr{R}_f(X)}\to\|\mu\|_{\mathscr{R}_f(X)}$（见\cref{remark:ch05-窄收敛范数}）。显然这蕴含 $M$ 有界。

为证 $M$ 胎紧，初步注意 $X$ 为 Polish 空间，故存在相容度量 $d$ 使 $(X,d)$ 完备。固定 $\epsilon>0$，设 $D\equiv\{x_k\}$ 为 $X$ 的可数稠密子集。对任意固定的 $p\in\mathbb{N}$，置 $E_{n,p}:=\bigcup_{k=1}^n B_{1/p}(x_k)$（其中 $B_{1/p}(x_k):=\{x\in X\mid d(x_k,x)<1/p\}$；$n\in\mathbb{N}$）。显然每个 $E_{n,p}$ 为开集，对所有 $n\in\mathbb{N}$ 有 $E_{n,p}\subseteq E_{n+1,p}$，且 $\bigcup_{n=1}^\infty E_{n,p}=X$。我们证明如下断言。

\textbf{断言。} 设 $\{A_n\}\subseteq\mathcal{P}(X)$ 为递增开集序列使 $\bigcup_{n=1}^\infty A_n=X$。则对每个 $\epsilon>0$，存在 $\bar{n}\in\mathbb{N}$ 使对所有 $\mu\in M$ 有 $\mu(A_{\bar{n}}^c)<\epsilon$。

由断言，对任意 $p\in\mathbb{N}$，存在 $n_p\in\mathbb{N}$ 使对所有 $\mu\in M$ 有 $\mu((E_{n_p,p})^c)<\epsilon/2^p$，特别地
\begin{align}
\mu((E_{n_p,p})^c)<\frac{\epsilon}{2^p}
\quad\text{对所有 }\mu\in M.\label{align:ch05-5.24}
\end{align}
置
\begin{align*}
K:=\bigcap_{p=1}^\infty E_{n_p,p}\subseteq\bigcap_{p=1}^\infty\bigcup_{k=1}^{n_p}B_{1/p}(x_k).
\end{align*}
由\eqref{align:ch05-5.24}，对所有 $\mu\in M$ 有 $\mu(K^c)<\epsilon$，故若我们证明 $K$ 紧，结论即得。为此注意，对每个 $p\in\mathbb{N}$，集合 $K$ 包含在有限并 $\bigcup_{k=1}^{n_p}B_{1/p}(x_k)$ 中，且 $\operatorname{diam}B_{1/p}(x_k)=2/p$（$p\in\mathbb{N}$）。由对角论证，每个序列 $\{x'_j\}\subseteq K$ 含 Cauchy 子列 $\{x'_{j_l}\}\subseteq\{x'_j\}$。由于 $K$ 闭且 $X$ 完备，存在 $x'\in K$ 使当 $l\to\infty$ 时 $d(x'_{j_l},x')\to0$。故 $K$ 序列紧，从而紧。结论得证。

尚需证明断言。反设存在递增开集序列 $\{A_n\}$ 使 $\bigcup_{n=1}^\infty A_n=X$，及 $\epsilon>0$ 具有如下性质：对所有 $n\in\mathbb{N}$，存在 $\mu_n\in M$ 使 $\mu_n(A_n^c)\geqslant\epsilon$。由于 $M$ 在窄拓扑中相对序列紧，存在 $\{\mu_{n_k}\}\subseteq\{\mu_n\}$ 与 $\mu\in\mathscr{R}_f(X)$ 使 $\mu_{n_k}\stackrel{n}{\rightharpoonup}\mu$。于是由\cref{theorem:ch05-portmanteau}与 $\{A_n\}$ 的单调性，对每个固定的 $n\in\mathbb{N}$，
\begin{align*}
\mu(A_n^c)\geqslant\limsup_{k\to\infty}\mu_{n_k}(A_n^c)
\geqslant\limsup_{k\to\infty}\mu_{n_k}(A_{n_k}^c)\geqslant\epsilon.
\end{align*}
然而由于 $\mu$ 有限，我们得 $\lim_{n\to\infty}\mu(A_n^c)=0$，矛盾。故断言成立，证明完成。
\end{proof}

\begin{proposition}\label{proposition:ch05-有界胎紧紧}
设 $X$ 为具可数基的局部紧 Hausdorff 空间，$M\subseteq\mathscr{R}^+_f(X)$ 有界且胎紧。则 $M$ 在窄拓扑中相对序列紧。
\end{proposition}

\begin{proof}
由假设，
\begin{align}
\sup_{\mu\in M}\|\mu\|_{\mathscr{R}_f(X)}
=\sup_{\mu\in M}\mu(X)=:C<\infty.\label{align:ch05-5.25}
\end{align}
回忆由\cref{proposition:appA-紧生成}，空间 $X$ 为 $\sigma$-紧。我们先在 $X$ 紧的假设下证明结果。设 $D\equiv\{f_k\}$ 为 $C_b(X)$ 的可数稠密子集（由于 $X$ 紧，$C_b(X)$ 可分；见\cref{proposition:appA-紧可分}）。则
\begin{align*}
\sup_{\mu\in M}\left|\int_X f_1\,\mathrm{d}\mu\right|
\leqslant\|f_1\|_\infty\sup_{\mu\in M}\|\mu\|_{\mathscr{R}_f(X)}
=\|f_1\|_\infty C<\infty,
\end{align*}
故集合 $\{\int_X f_1\,\mathrm{d}\mu\mid\mu\in M\}\subseteq\mathbb{R}$ 有界。于是存在序列 $\{\mu_j^1\}\subseteq M$ 与 $a_1\in\mathbb{R}$ 使
\begin{align}
\lim_{j\to\infty}\int_X f_1\,\mathrm{d}\mu_j^1=a_1.\label{align:ch05-5.26}
\end{align}
由于
\begin{align*}
\sup_{j\in\mathbb{N}}\left|\int_X f_2\,\mathrm{d}\mu_j^1\right|
\leqslant\|f_2\|_\infty\sup_{j\in\mathbb{N}}\|\mu_j^1\|
\leqslant\|f_2\|_\infty C<\infty,
\end{align*}
集合 $\{\int_X f_2\,\mathrm{d}\mu_j^1\mid j\in\mathbb{N}\}\subseteq\mathbb{R}$ 也有界。于是存在子列 $\{\mu_j^2\}\subseteq\{\mu_j^1\}$ 与 $a_2\in\mathbb{R}$ 使
\begin{align*}
\lim_{j\to\infty}\int_X f_2\,\mathrm{d}\mu_j^2=a_2.
\end{align*}
显然由\eqref{align:ch05-5.26}，我们也有 $\lim_{j\to\infty}\int_X f_1\,\mathrm{d}\mu_j^2=a_1$。迭代此论证，对每个 $k\in\mathbb{N}$，$k\geqslant3$，存在子列 $\{\mu_j^k\}\subseteq\{\mu_j^{k-1}\}$ 与 $a_k\in\mathbb{R}$ 使
\begin{align*}
\lim_{j\to\infty}\int_X f_l\,\mathrm{d}\mu_j^k=a_l
\quad\text{对所有 }l=1,\dots,k.
\end{align*}
由对角论证，序列 $\{\mu_j\}\subseteq M$（其中 $\mu_j:=\mu_j^j$，$j\in\mathbb{N}$）满足
\begin{align}
\lim_{j\to\infty}\int_X f_k\,\mathrm{d}\mu_j=a_k
\quad\text{对所有 }k\in\mathbb{N}.\label{align:ch05-5.27}
\end{align}
现考虑映射 $F_0:D\to\mathbb{R}$，$F_0(f_k):=a_k$（$k\in\mathbb{N}$）。由\eqref{align:ch05-5.27}，$F_0$ 线性、正且连续，因为
\begin{align*}
|F_0(f_k)|=|a_k|
\leqslant\lim_{j\to\infty}\left|\int_X f_k\,\mathrm{d}\mu_j\right|
\leqslant\|f_k\|_\infty C
\quad\text{对所有 }k\in\mathbb{N}.
\end{align*}
由于 $D$ 在 $C_b(X)$ 中稠密，$F_0$ 可唯一延拓为 $(C_b(X))^*$ 上的泛函 $F$，它是正的且满足 $\|F\|\leqslant C$。我们证明如下断言。

\textbf{断言。} 对任意 $\epsilon>0$，存在紧子集 $K\subseteq X$ 使对所有 $f\in C_b(X)$ 满足 $f(X)\subseteq[0,1]$ 且 $f|_K=0$，有 $\langle F,f\rangle<\epsilon$。

依靠断言，我们可在 $X$ 紧时完成证明。事实上，由\cref{theorem:ch02-Riesz表示}，存在唯一 $\mu\in\mathscr{R}^+_f(X)$ 使
\begin{align}
\langle F,f\rangle=\int_X f\,\mathrm{d}\mu
\quad\text{对所有 }f\in C_b(X)\label{align:ch05-5.28}
\end{align}
（见\eqref{align:ch02-2.49}）；此外由\cref{theorem:ch02-Riesz范数}(ii)，$\|\mu\|_{\mathscr{R}_f(X)}=\|F\|\leqslant C$。设 $g\in C_b(X)$，$\{f_l\}\subseteq D$ 为使 $\lim_{l\to\infty}\|f_l-g\|_\infty=0$ 的序列。固定 $\epsilon>0$，取 $l_0\in\mathbb{N}$ 充分大使 $\|f_{l_0}-g\|_\infty<\epsilon/(4C)$。由\eqref{align:ch05-5.27}–\eqref{align:ch05-5.28}，
\begin{align*}
\langle F,f_{l_0}\rangle=F_0(f_{l_0})=a_{l_0}
=\int_X f_{l_0}\,\mathrm{d}\mu
=\lim_{j\to\infty}\int_X f_{l_0}\,\mathrm{d}\mu_j,
\end{align*}
故存在 $j_0\in\mathbb{N}$ 使
\begin{align*}
\left|\int_X f_{l_0}\,\mathrm{d}\mu_j-\int_X f_{l_0}\,\mathrm{d}\mu\right|<\frac{\epsilon}{2}
\quad\text{对所有 }j>j_0.
\end{align*}
则对所有 $j>j_0$，
\begin{align*}
\left|\int_X g\,\mathrm{d}\mu_j-\int_X g\,\mathrm{d}\mu\right|
\leqslant\left|\int_X(g-f_{l_0})\,\mathrm{d}\mu_j\right|
+\left|\int_X(g-f_{l_0})\,\mathrm{d}\mu\right|
+\left|\int_X f_{l_0}\,\mathrm{d}\mu_j-\int_X f_{l_0}\,\mathrm{d}\mu\right|\\
\leqslant\|f_{l_0}-g\|_\infty\left(\sup_{j\in\mathbb{N}}\|\mu_j\|+\|\mu\|_{\mathscr{R}_f(X)}\right)+\frac{\epsilon}{2}
<2C\cdot\frac{\epsilon}{4C}+\frac{\epsilon}{2}=\epsilon.
\end{align*}
综上，我们证明了存在序列 $\{\mu_j\}\subseteq M$ 与 $\mu\in\mathscr{R}^+_f(X)$ 使 $\mu_j\stackrel{n}{\rightharpoonup}\mu$，结论得证。

现证明上述断言。由于由假设 $M$ 胎紧，对任意 $\epsilon>0$，存在紧子集 $K\subseteq X$ 使 $\sup_{\mu\in M}\mu(K^c)<\epsilon/6$。设 $f\in C_b(X)$ 满足 $f(X)\subseteq[0,1]$ 且 $f|_K=0$。由于 $D$ 在 $C_b(X)$ 中稠密，存在序列 $\{f_l\}\subseteq D$（对所有 $l\in\mathbb{N}$，$f_l\geqslant0$）使 $\lim_{l\to\infty}\|f_l-f\|_\infty=0$；因而特别地，对所有充分大的 $l\in\mathbb{N}$，$\|f_l\|_\infty<1+\epsilon/6$，$\|f_l|_K\|_\infty<\epsilon/(6C)$。

则对任意充分大的 $l\in\mathbb{N}$，
\begin{align*}
\langle F,f_l\rangle=a_l
=\lim_{j\to\infty}\int_X f_l\,\mathrm{d}\mu_j
\leqslant\limsup_{j\to\infty}\int_K f_l\,\mathrm{d}\mu_j
+\limsup_{j\to\infty}\int_{K^c}f_l\,\mathrm{d}\mu_j\\
\leqslant\frac{\epsilon}{6C}\sup_{j\in\mathbb{N}}\mu_j(X)
+\left(1+\frac{\epsilon}{6}\right)\sup_{j\in\mathbb{N}}\mu_j(K^c)
<\frac{\epsilon}{6}+\left(1+\frac{\epsilon}{6}\right)\frac{\epsilon}{6}
<\frac{\epsilon}{2}.
\end{align*}
由于 $\|f_l-f\|_\infty\to0$ 且 $F$ 连续，由该不等式得
\begin{align*}
\langle F,f\rangle=\lim_{l\to\infty}\langle F,f_l\rangle
\leqslant\frac{\epsilon}{2}<\epsilon,
\end{align*}
故断言成立。$X$ 紧时的证明完成。

一般情形，设 $\{K_l\}$ 为 $X$ 的递增紧子集序列使 $\bigcup_{l=1}^\infty K_l=X$。由\eqref{align:ch05-5.25}，对每个 $l\in\mathbb{N}$，
\begin{align*}
\sup_{\mu\in M}\|\mu\|_{\mathscr{R}_f(K_l)}
\leqslant\sup_{\mu\in M}\|\mu\|_{\mathscr{R}_f(X)}=C.
\end{align*}
由上述考虑，对每个 $l\in\mathbb{N}$，存在序列 $\{\mu_j^l\}\subseteq M$ 与 $\mu^l\in\mathscr{R}^+_f(X)$ 使当 $j\to\infty$ 时 $\mu_j^l\stackrel{n}{\rightharpoonup}\mu^l$。于是由对角论证结论得证。
\end{proof}

由\cref{theorem:ch05-prokhorov}，拥有胎紧性判据是有用的。我们给出如下定义。

\begin{definition}\label{definition:ch05-下紧}
设 $X$ 为局部紧 Hausdorff 空间。函数 $\varphi:X\to(-\infty,\infty]$ 称为\textbf{下紧的}（inf-compact），若对每个 $r\in\mathbb{R}$，集合 $\{\varphi\leqslant r\}:=\{x\in X\mid\varphi(x)\leqslant r\}$ 为紧集。
\end{definition}

\begin{example}\label{example:ch05-下紧}
若 $X=\mathbb{R}$，则每个连续函数 $\varphi:\mathbb{R}\to[0,\infty)$ 满足当 $|x|\to\infty$ 时 $\varphi(x)\to\infty$ 都是下紧的。事实上，对每个 $r\in\mathbb{R}$，集合 $\{\varphi\leqslant r\}$ 闭且有界，故紧。
\end{example}

\cref{definition:ch05-下紧}导出 $\mathscr{R}^+_f(X)$ 胎紧子集的如下刻画。

\begin{proposition}\label{proposition:ch05-胎紧判据}
设 $X$ 为具可数基的局部紧 Hausdorff 空间。则下列命题等价：
\begin{enumerate}[(i)]
\item\label{item:ch05-tight-i} $M\subseteq\mathscr{R}^+_f(X)$ 胎紧；
\item\label{item:ch05-tight-ii} 存在下紧函数 $\varphi:X\to[0,\infty)$ 使
\begin{align}
\sup_{\mu\in M}\int_X\varphi\,\mathrm{d}\mu<\infty.\label{align:ch05-5.29}
\end{align}
\end{enumerate}
\end{proposition}

\begin{proof}
(i)$\Rightarrow$(ii)。由\cref{proposition:appA-紧生成}，存在 $X$ 的递增紧子集序列 $\{K_n\}$ 使对所有 $\mu\in M$ 有 $\mu(K_n^c)<1/2^n$。置 $\varphi:=\sum_{n=1}^\infty\chi_{K_n^c}$。则对每个 $\mu\in M$，
\begin{align*}
\sup_{\mu\in M}\int_X\varphi\,\mathrm{d}\mu
\leqslant\sum_{n=1}^\infty\frac{1}{2^n}<\infty,
\end{align*}
故不等式\eqref{align:ch05-5.29}成立。

为证 $\varphi$ 下紧，固定 $x\in X$。对每个 $n\in\mathbb{N}$，我们有 (a) $x\in K_n$ 或 (b) $x\in K_n^c$。由于每个 $n$ 有 $K_n^c\supseteq K_{n+1}^c$，
\begin{itemize}
\item 在情形 (a)，$\varphi(x)\leqslant n-1$；
\item 在情形 (b)，若 $x\in K_n^c\setminus K_{n+1}^c=K_n^c\cap K_{n+1}$ 则 $\varphi(x)=n$，若 $x\in K_{n+1}^c$ 则 $\varphi(x)\geqslant n+1$。
\end{itemize}
由此 $\{\varphi\leqslant n\}=K_{n+1}$，故显然对所有 $r>0$ 有 $\{\varphi\leqslant r\}=K_{[r]+1}$，其中 $[r]$ 表示小于 $r$ 的最大整数。由于 $K_{[r]+1}$ 紧，断言得证。

(ii)$\Rightarrow$(i)。由假设存在 $\varphi:X\to[0,\infty)$ 使对每个 $r>0$ 集合 $K_r\equiv\{\varphi\leqslant r\}$ 紧，且\eqref{align:ch05-5.29}成立。则由 Chebyshev 不等式，对每个 $\mu\in M$，
\begin{align*}
r\mu(K_r^c)<\int_{K_r^c}\varphi\,\mathrm{d}\mu
\leqslant\int_X\varphi\,\mathrm{d}\mu
\leqslant\sup_{\mu\in M}\int_X\varphi\,\mathrm{d}\mu,
\end{align*}
由 $r$ 的任意性结论得证。
\end{proof}

\begin{remark}\label{remark:ch05-径向下紧}
由上述证明，若 $X=\mathbb{R}$，则可取 $K_n=B_n(0)$，即中心在 $0$、半径 $n\in\mathbb{N}$ 的闭球。在此情形，$\varphi$ 可写为
\begin{align*}
\varphi(x):=\sum_{n=0}^\infty n\,\chi_{(n,n+1]}(|x|)\qquad(x\in X);
\end{align*}
注意它径向、非减且在无穷远发散。
\end{remark}

\section{参数化测度与离散分解}

设 $X$ 与 $Y$ 为局部紧 Hausdorff 空间，$\mathscr{B}(X)$ 与 $\mathscr{B}(Y)$ 为相应 Borel $\sigma$-代数。回忆由\cref{theorem:ch02-Rf对偶}，Banach 空间 $\mathscr{R}_f(X)$ 是 $C_0(X)$ 的对偶空间。于是我们可给出如下定义。

\begin{definition}\label{definition:ch05-参数测度}
设 $X$ 与 $Y$ 为局部紧 Hausdorff 空间。从 $X$ 到 $\mathscr{R}_f(Y)$ 的弱$^*$可测映射称为 $Y$ 上的\textbf{参数化测度}。

对 $Y$ 上的参数化测度，我们使用记号 $\{\nu_x\}_{x\in X}$，其中对每个 $x\in X$ 有 $\nu_x\in\mathscr{R}_f(Y)$。
\end{definition}

\begin{remark}\label{remark:ch05-参数测度等价}
我们指出\cref{definition:ch05-参数测度}的两种等价表述（例如见文献~\cite{ref05}）。首先，由\cref{definition:ch04-弱星可测}(ii)与\cref{theorem:ch02-Rf对偶}，族 $\{\nu_x\}_{x\in X}\subseteq\mathscr{R}_f(Y)$ 是 $Y$ 上的参数化测度，当且仅当
\begin{align}
x\mapsto\langle\nu_x,g\rangle=\int_Y g(y)\,\mathrm{d}\nu_x(y)
\text{ 对每个 }g\in C_0(Y)\text{ 为 }\mathscr{B}(X)\text{-可测}.\label{align:ch05-5.30}
\end{align}
其次，条件\eqref{align:ch05-5.30}可重述如下：
\begin{align}
x\mapsto\nu_x(F)\text{ 对每个 }F\in\mathscr{B}(Y)\text{ 为 }\mathscr{B}(X)\text{-可测}.\label{align:ch05-5.31}
\end{align}
事实上，由\eqref{align:ch05-5.31}，映射 $x\mapsto\int_Y\chi_F(y)\,\mathrm{d}\nu_x(y)$ 对每个 $F\in\mathscr{B}(Y)$ 为 $\mathscr{B}(X)$-可测，故将 $\chi_F$ 替换为任意简单函数 $s\in\mathcal{S}(Y)$ 后同样成立。另一方面，对每个 $g\in C_0(Y)$，存在序列 $\{s_n\}\subseteq\mathcal{S}(Y)$ 使当 $n\to\infty$ 时 $\|s_n-g\|_\infty\to0$（见\cref{theorem:ch02-简单函数逼近}(ii)）。故对所有 $x\in X$，
\begin{align*}
\lim_{n\to\infty}\int_Y[s_n(y)-g(y)]\,\mathrm{d}\nu_x(y)=0,
\end{align*}
于是由\cref{corollary:ch02-控制收敛}(ii)得\eqref{align:ch05-5.30}。

反之，设\eqref{align:ch05-5.30}成立，$F\in\mathscr{B}(Y)$。由\cref{proposition:ch02-连续逼近}（另见\cref{remark:ch02-连续逼近注}），存在序列 $\{\zeta_m\}\subseteq C_c(Y)$ 使 $\zeta_m\to\chi_F$ 在 $Y$ 中 $\mu$-几乎处处成立，且对所有 $m\in\mathbb{N}$ 有 $\|\zeta_m\|_\infty\leqslant1$。然后在\eqref{align:ch05-5.30}中取 $g=\zeta_m$，令 $m\to\infty$，并用控制收敛定理，即得\eqref{align:ch05-5.31}。
\end{remark}

\begin{remark}\label{remark:ch05-变差可测}
若 $\{\nu_x\}_{x\in X}$ 为 $Y$ 上的参数化测度，则实值函数 $x\mapsto|\nu_x|(F)$ 对每个 $F\in\mathscr{B}(Y)$ 也为 $\mathscr{B}(X)$-可测。事实上，由\cref{corollary:ch02-控制收敛}(i)，函数 $x\mapsto\sum_{i=1}^n|\nu_x(F_i)|$ 对 $F$ 的任意划分 $\{F_1,\dots,F_n\}\subseteq\mathscr{B}(Y)$ 为 $\mathscr{B}(X)$-可测，故由\cref{proposition:ch01-全变差}(见\eqref{align:ch01-全变差公式})与\cref{corollary:ch02-控制收敛}(ii)断言得证。
\end{remark}

设 $X$ 与 $Y$ 为具可数基的局部紧 Hausdorff 空间。则拓扑积 $X\times Y$ 为具可数基的局部紧 Hausdorff 空间，且由\cref{theorem:ch01-Borel积}
\begin{align*}
\mathscr{B}(X\times Y)=\mathscr{B}(X)\otimes\mathscr{B}(Y).
\end{align*}
由\cref{lemma:ch02-投影可测}，投影
\begin{align*}
p_X:X\times Y\to X,\qquad p_X(x,y):=x
\end{align*}
为 $(\mathscr{B}(X\times Y),\mathscr{B}(X))$-可测，投影
\begin{align*}
p_Y:X\times Y\to Y,\qquad p_Y(x,y):=y
\end{align*}
为 $(\mathscr{B}(X\times Y),\mathscr{B}(Y))$-可测。故对任意 $\mu\in\mathcal{M}_f(X\times Y)$，像测度（见\eqref{align:ch02-2.28a}）
\begin{align}
\mu_{p_X}:\mathscr{B}(X)\to\mathbb{R},\quad
\mu_{p_X}(E):=\mu(p_X^{-1}(E))=\mu(E\times Y)
\quad\text{对所有 }E\in\mathscr{B}(X),\label{align:ch05-5.32}\\
\mu_{p_Y}:\mathscr{B}(Y)\to\mathbb{R},\quad
\mu_{p_Y}(F):=\mu(p_Y^{-1}(F))=\mu(X\times F)
\quad\text{对所有 }F\in\mathscr{B}(Y),\label{align:ch05-5.33}
\end{align}
是良定义的。若 $\mu\in\mathscr{R}_f(X\times Y)$，则由\cref{proposition:ch01-Radon投影}，$\mu_{p_X}\in\mathscr{R}_f(X)$，$\mu_{p_Y}\in\mathscr{R}_f(Y)$。

\begin{theorem}[离散分解]\label{theorem:ch05-离散分解}
设 $X$ 与 $Y$ 为具可数基的局部紧 Hausdorff 空间。设 $\nu\in\mathscr{R}_f(X\times Y)$，$|\nu|_{p_X}\in\mathscr{R}^+_f(X)$ 为其全变差 $|\nu|$ 到空间 $X$ 的投影。则存在 $Y$ 上的参数化测度 $\nu\equiv\{\nu_x\}_{x\in X}$ 具有下列性质：
\begin{enumerate}[(i)]
\item\label{item:ch05-dis-i} $|\nu_x|(Y)=1$ 对 $|\nu|_{p_X}$-几乎所有 $x\in X$；
\item\label{item:ch05-dis-ii} 对每个 $h\in L^1(X\times Y,\mathscr{B}(X\times Y),|\nu|)$，
\begin{align}
h(x,\cdot)\in L^1(Y,\mathscr{B}(Y),|\nu_x|)
\quad\text{对 }|\nu|_{p_X}\text{-几乎所有 }x\in X,\label{align:ch05-5.34}\\
\text{映射 }x\mapsto\int_Y h(x,y)\,\mathrm{d}\nu_x(y)
\text{ 属于 }L^1(X,\mathscr{B}(X),|\nu|_{p_X}),\label{align:ch05-5.35}
\end{align}
且
\begin{align}
\int_{X\times Y}h(x,y)\,\mathrm{d}\nu(x,y)
=\int_X\mathrm{d}|\nu|_{p_X}(x)\int_Y h(x,y)\,\mathrm{d}\nu_x(y).\label{align:ch05-5.36}
\end{align}
\end{enumerate}
此外，参数化测度 $\{\nu_x\}_{x\in X}$ 在如下意义下唯一：设 $\{\nu'_x\}_{x\in X}$ 为 $Y$ 上的参数化测度使
\begin{align}
\text{映射 }x\mapsto\int_Y h(x,y)\,\mathrm{d}\nu'_x(y)
\text{ 属于 }L^1(X,\mathscr{B}(X),|\nu|_{p_X})\label{align:ch05-5.37}
\end{align}
且
\begin{align}
\int_{X\times Y}h(x,y)\,\mathrm{d}\nu(x,y)
=\int_X\mathrm{d}|\nu|_{p_X}(x)\int_Y h(x,y)\,\mathrm{d}\nu'_x(y)\label{align:ch05-5.38}
\end{align}
对每个有界 $\mathscr{B}(X\times Y)$-可测、紧支撑的 $h:X\times Y\to\mathbb{R}$ 成立。则 $\nu'_x=\nu_x$ 对 $|\nu|_{p_X}$-几乎所有 $x\in X$。
\end{theorem}

\begin{definition}\label{definition:ch05-离散分解定义}
参数化测度 $\{\nu_x\}_{x\in X}$ 称为测度 $\nu$ 的\textbf{离散分解}（disintegration）。
\end{definition}

\begin{remark}\label{remark:ch05-离散分解注}
(i) 设 $G\in\mathscr{B}(X\times Y)$，$G_x:=\{y\in Y\mid(x,y)\in G\}\in\mathscr{B}(Y)$ 为其 $x$-截口（见\cref{proposition:ch01-截口}）。在\eqref{align:ch05-5.36}中取 $h=\chi_G$，并注意 $\chi_G(x,\cdot)=\chi_{G_x}$（见\cref{remark:ch02-截口}），得
\begin{align}
\nu(G)=\int_X\nu_x(G_x)\,\mathrm{d}|\nu|_{p_X}(x)
\quad\text{对所有 }G\in\mathscr{B}(X\times Y).\label{align:ch05-5.39}
\end{align}
若 $\nu$ 正，则由\eqref{align:ch05-5.39}与 $G$ 的任意性，$\nu_x$ 对 $|\nu|_{p_X}$-几乎所有 $x\in X$ 为正。

(ii) 设 $\nu_{p_Y}\in\mathscr{R}_f(Y)$ 为 $\nu$ 到 $Y$ 的投影（见\eqref{align:ch05-5.33}）。对任意 $F\in\mathscr{B}(Y)$，在\eqref{align:ch05-5.39}中取 $G=X\times F$（从而对每个 $x\in X$ 有 $G_x=F$），得
\begin{align}
\nu_{p_Y}(F)=\nu(X\times F)=\int_X\nu_x(F)\,\mathrm{d}|\nu|_{p_X}(x).\label{align:ch05-5.40}
\end{align}
\end{remark}

\begin{proof}[证明\cref{theorem:ch05-离散分解}]
对任意 $g\in C_0(Y)$，$g=g(y)$，置
\begin{align*}
\mu_g:\mathscr{B}(X)\to[0,\infty),\quad
\mu_g(E):=\int_{E\times Y}g\,\mathrm{d}\nu
\quad\text{对所有 }E\in\mathscr{B}(X).
\end{align*}
显然 $\mu_g\in\mathscr{R}_f(X)$，且
\begin{align}
|\mu_g(E)|\leqslant\|g\|_\infty|\nu|(E\times Y)
=\|g\|_\infty|\nu|_{p_X}(E)
\quad\text{对所有 }E\in\mathscr{B}(X).\label{align:ch05-5.42}
\end{align}
由不等式\eqref{align:ch05-5.42}，$\mu_g\ll|\nu|_{p_X}$，故由 Radon–Nikodým 定理存在 $(\mu_g)_r\in L^1(X,\mathscr{B}(X),|\nu|_{p_X})$ 使
\begin{align}
\mu_g(E)=\int_E(\mu_g)_r\,\mathrm{d}|\nu|_{p_X}
\quad\text{对所有 }E\in\mathscr{B}(X).\label{align:ch05-5.43}
\end{align}
此外，由\eqref{align:ch05-5.42}与\cref{proposition:ch04-有界密度}(见\eqref{align:ch04-4.43b})，$(\mu_g)_r\in L^\infty(X,\mathscr{B}(X),|\nu|_{p_X})$，且 $\|(\mu_g)_r\|_\infty\leqslant\|g\|_\infty$。显然，映射 $g\mapsto\mu_g$（$g\in C_0(Y)$，$\mu_g$ 由\eqref{align:ch05-5.42}定义）是线性的。设 $D\subseteq C_0(Y)$ 可数且稠密（注意空间 $C_0(Y)$ 可分；见附录 A,\cref{paragraph:appA-可分}）。则由\eqref{align:ch05-5.43}，存在 $|\nu|_{p_X}$-零测集 $N\in\mathscr{B}(X)$ 使
\begin{align}
(\mu_{\alpha g+\beta g'})_r(x)=\alpha(\mu_g)_r(x)+\beta(\mu_{g'})_r(x)
\quad\text{对所有 }x\in N^c,\ \alpha,\beta\in\mathbb{R},\ g,g'\in D.\label{align:ch05-5.44}
\end{align}
置
\begin{align*}
T_x:D\to\mathbb{R},\qquad T_xg:=(\mu_g)_r(x)
\quad\text{对所有 }x\in N^c\text{ 与 }g\in D.
\end{align*}

由\eqref{align:ch05-5.44}，$T_x$ 线性。此外，
\begin{align}
|T_xg|=|(\mu_g)_r(x)|\leqslant\|(\mu_g)_r\|_\infty\leqslant\|g\|_\infty
\qquad(x\in N^c),\label{align:ch05-5.45}
\end{align}
故 $T_x$ 有界。由于 $D$ 在 $C_0(Y)$ 中稠密，$T_x$ 可唯一延拓（仍记为 $T_x$）到整个空间 $C_0(Y)$。因而对所有 $x\in N^c$，$T_x\in(C_0(Y))^*$，故由\cref{theorem:ch02-Rf对偶}存在 $\nu_x\in\mathscr{R}_f(Y)$ 使
\begin{align}
T_xg=(\mu_g)_r(x)=\int_Y g(y)\,\mathrm{d}\nu_x(y)
\quad\text{对所有 }g\in C_0(Y)\ (x\in N^c)\label{align:ch05-5.46}
\end{align}
且
\begin{align}
\|\nu_x\|=|\nu_x|(Y)=\|T_x\|\leqslant1\label{align:ch05-5.47}
\end{align}
（见\eqref{align:ch05-5.45}）。我们可通过对任意 $x\in N$ 置 $\nu_x:=\delta_{\bar{y}}$（其中 $\bar{y}\in Y$ 任意固定）来补全 $\nu_x$ 在 $X$ 上的定义。

我们证明函数 $x\mapsto\nu_x(F)$ 对每个 $F\in\mathscr{B}(Y)$ 为 $\mathscr{B}(X)$-可测，从而族 $\{\nu_x\}_{x\in X}$ 为 $Y$ 上的参数化测度（见\cref{remark:ch05-参数测度等价}）。事实上，由定义对所有 $x\in N^c$ 有 $T_xg=(\mu_g)_r(x)$；由于 $(\mu_g)_r\in L^\infty(X,\mathscr{B}(X),|\nu|_{p_X})$，因而为 $\mathscr{B}(X)$-可测，映射 $x\mapsto T_xg$ 对任意 $g\in C_0(Y)$ 为 $(\mathscr{B}(X)\cap N^c)$-可测。现注意对每个 $F\in\mathscr{B}(Y)$，存在序列 $\{g_n\}\subseteq C_0(Y)$ 使当 $n\to\infty$ 时 $g_n\to\chi_F$ 在 $L^1(Y,\mathscr{B}(Y),\nu_x)$ 中（见\cref{proposition:ch02-连续逼近}与\cref{remark:ch02-连续逼近注}）。则
\begin{align*}
T_xg_n=\int_Y g_n(y)\,\mathrm{d}\nu_x(y)\to\nu_x(F)
\quad(n\to\infty),
\end{align*}
故由\cref{corollary:ch02-控制收敛}(ii)，映射 $x\mapsto\nu_x(F)$ 为 $\mathscr{B}(X)\cap N^c$-可测。另一方面，显然映射 $x\mapsto\nu_x(F)=\delta_{\bar{y}}(F)$（$x\in N$）对所有 $F\in\mathscr{B}(Y)$ 为 $(\mathscr{B}(X)\cap N)$-可测。于是由\cref{proposition:ch02-可测拼接}断言得证。

鉴于上述注记，显然每个形如 $h(x,y):=s(x)g(y)$ 的映射（$s\in\mathcal{S}(X)$，$g\in C_0(Y)$）满足\eqref{align:ch05-5.34}–\eqref{align:ch05-5.35}（见\cref{remark:ch05-参数测度等价}）。此外，等式\eqref{align:ch05-5.36}在此 $h(x,y)$ 的选取下也成立。事实上，由\eqref{align:ch05-5.42}、\eqref{align:ch05-5.43}与\eqref{align:ch05-5.46}，对任意 $g\in C_0(Y)$ 与 $E\in\mathscr{B}(X)$，
\begin{align*}
\mu_g(E)=\int_{X\times Y}\chi_E(x)g(y)\,\mathrm{d}\nu
=\int_E(\mu_g)_r(x)\,\mathrm{d}|\nu|_{p_X}(x)
=\int_X\mathrm{d}|\nu|_{p_X}(x)\int_Y\chi_E(x)g(y)\,\mathrm{d}\nu_x(y),
\end{align*}
由此
\begin{align*}
\int_{X\times Y}s(x)g(y)\,\mathrm{d}\nu
=\int_X\mathrm{d}|\nu|_{p_X}(x)\int_Y s(x)g(y)\,\mathrm{d}\nu_x(y)
\end{align*}
对所有 $s\in\mathcal{S}(X)$ 与 $g\in C_0(Y)$ 成立。由\cref{lemma:ch02-简单函数}与控制收敛定理，上述等式对所有 $s\in\mathcal{S}(X)$ 与 $g\in\mathcal{S}(Y)$ 也成立。由于集合 $\{s(x)g(y)\mid s\in\mathcal{S}(X),\ g\in\mathcal{S}(Y)\}$ 在 $L^1(X\times Y,\mathscr{B}(X\times Y),\nu)$ 中稠密，断言 (ii) 得证。

由\cref{remark:ch05-变差可测}、\cref{proposition:ch01-全变差}与\cref{corollary:ch02-控制收敛}(ii)，对每个 $F\in\mathscr{B}(Y)$，函数 $x\mapsto|\nu_x|(F)$ 为 $\mathscr{B}(X)$-可测。于是由\eqref{align:ch05-5.39}，
\begin{align*}
|\nu(G)|\leqslant\int_X|\nu_x|(G_x)\,\mathrm{d}|\nu|_{p_X}(x)
\quad\text{对所有 }G\in\mathscr{B}(X\times Y),
\end{align*}
故由 $G$ 的任意性与 $|\nu|$ 的定义，
\begin{align*}
|\nu|(G)\leqslant\int_X|\nu_x|(G_x)\,\mathrm{d}|\nu|_{p_X}(x).
\end{align*}
在此不等式中取 $G=X\times Y$ 并用\eqref{align:ch05-5.47}，得
\begin{align*}
|\nu|(X\times Y)\leqslant\int_X|\nu_x|(Y)\,\mathrm{d}|\nu|_{p_X}(x)
\leqslant|\nu|(X\times Y)
\end{align*}
（见\eqref{align:ch05-5.32}），故断言 (i) 得证。

最后，设 $\{\nu'_x\}_{x\in X}$ 为 $Y$ 上的参数化测度，对所有有界 $\mathscr{B}(X\times Y)$-可测紧支撑 $h$ 满足\eqref{align:ch05-5.37}–\eqref{align:ch05-5.38}。由\eqref{align:ch05-5.36}与\eqref{align:ch05-5.38}得
\begin{align*}
\int_X\mathrm{d}|\nu|_{p_X}(x)\int_Y h(x,y)\,\mathrm{d}\nu_x(y)
=\int_X\mathrm{d}|\nu|_{p_X}(x)\int_Y h(x,y)\,\mathrm{d}\nu'_x(y)
\end{align*}
对所有上述 $h$ 成立，故显然 $\nu'_x=\nu_x$ 对 $|\nu|_{p_X}$-几乎所有 $x\in X$。证明完成。
\end{proof}

\begin{remark}\label{remark:ch05-离散分解绝对连续}
设 $X,Y$ 如\cref{theorem:ch05-离散分解}，$\nu\in\mathscr{R}^+_f(X\times Y)$，$\{\nu_x\}_{x\in X}$ 为其离散分解。由等式\eqref{align:ch05-5.40}，对每个使 $\nu_{p_Y}(F)=0$ 的 $F\in\mathscr{B}(Y)$，存在 $\nu_{p_X}$-零测集 $N\subseteq X$ 使对所有 $x\in N^c$ 有 $\nu_x(F)=0$。由于 $Y$ 具有可数基，可测空间 $(Y,\mathscr{B}(Y))$ 是可分的（见\cref{definition:ch01-可分空间}）。于是由标准论证，$\nu_{p_X}$-零测集 $N$ 的选取可独立于 $F\in\mathscr{B}(Y)$。由此 $\nu_x\ll\nu_{p_Y}$ 对 $\nu_{p_X}$-几乎所有 $x\in X$。
\end{remark}

\section{Young 测度再探}

设 $X$ 为局部紧 Hausdorff 空间，$\{u_k\}\subseteq L^1(X,\mathscr{B}(X),\mu)$ 弱收敛。已观察到（见\cref{example:ch02-弱收敛振荡}(i)），对非线性 $h$（使 $\{h\circ u_k\}\subseteq L^1(X,\mathscr{B}(X),\mu)$），$\{h\circ u_k\}$ 未必如此。为获得 $\{h\circ u_k\}$ 性态的信息，\cref{remark:ch05-函数转测度}建议我们处理与 $\{h\circ u_k\}$ 相伴的序列 $\{\nu_k\}\subseteq\mathscr{R}_f(X)$，如\eqref{align:ch05-5.5}。在此框架中，与 $\{u_k\}$ 相伴的 Young 测度序列 $\{\nu_{u_k}\}$ 起重要作用。

设 $\mu\in\mathscr{R}^+_f(X)$，$\nu\in\mathcal{Y}^+(X,\mathscr{B}(X),\mu;Y)\equiv\mathcal{Y}^+(X;Y)$ 为 Young 测度。由定义 $\nu_{p_X}=\mu$（见\cref{definition:ch02-Young测度}），且由于 $\mathcal{Y}^+(X;Y)\subseteq\mathscr{R}^+_f(X\times Y)$，由\cref{theorem:ch05-离散分解}得如下结果。

\begin{proposition}[Young 测度的离散分解]\label{proposition:ch05-离散分解}
设 $X$ 与 $Y$ 为具可数基的局部紧 Hausdorff 空间。设 $\mu\in\mathscr{R}^+_f(X)$，$\nu\in\mathcal{Y}^+(X;Y)$。则存在 $Y$ 上的参数化测度 $\nu\equiv\{\nu_x\}_{x\in X}$ 具有下列性质：
\begin{enumerate}[(i)]
\item\label{item:ch05-youdis-i} $\nu_x\in P(Y)\subseteq\mathscr{R}^+_f(Y)$ 对 $\mu$-几乎所有 $x\in X$；
\item\label{item:ch05-youdis-ii} 对每个 $h\in L^1(X\times Y,\mathscr{B}(X\times Y),\nu)$，
\begin{align}
h(x,\cdot)\in L^1(Y,\mathscr{B}(Y),\nu_x)
\quad\text{对 }\mu\text{-几乎所有 }x\in X,\label{align:ch05-5.48}\\
\text{映射 }x\mapsto\int_Y h(x,y)\,\mathrm{d}\nu_x(y)
\text{ 属于 }L^1(X,\mathscr{B}(X),\mu),\label{align:ch05-5.49}
\end{align}
且
\begin{align}
\int_{X\times Y}h(x,y)\,\mathrm{d}\nu(x,y)
=\int_X\mathrm{d}\mu(x)\int_Y h(x,y)\,\mathrm{d}\nu_x(y).\label{align:ch05-5.50}
\end{align}
\end{enumerate}
若 $\{\nu'_x\}_{x\in X}$ 为 $Y$ 上的参数化测度使
\begin{align}
\text{映射 }x\mapsto\int_Y h(x,y)\,\mathrm{d}\nu'_x(y)
\text{ 属于 }L^1(X,\mathscr{B}(X),\mu)\label{align:ch05-5.51}
\end{align}
且
\begin{align}
\int_{X\times Y}h(x,y)\,\mathrm{d}\nu(x,y)
=\int_X\mathrm{d}\mu(x)\int_Y h(x,y)\,\mathrm{d}\nu'_x(y)\label{align:ch05-5.52}
\end{align}
对每个有界 $\mathscr{B}(X\times Y)$-可测、紧支撑的 $h:X\times Y\to\mathbb{R}$ 成立，则 $\nu'_x=\nu_x$ 对 $\mu$-几乎所有 $x\in X$。
\end{proposition}

刻画与函数相伴的 Young 测度的离散分解是有趣的（见\cref{definition:ch02-Young测度}）。回忆与 $u$ 相伴的 Young 测度 $\nu_u$ 定义为
\begin{align*}
\nu_u(E\times F):=\mu(E\cap u^{-1}(F))
\quad\text{对所有 }E\in\mathscr{B}(X),\ F\in\mathscr{B}(Y).
\end{align*}

\begin{proposition}\label{proposition:ch05-Dirac测度}
设 $X$ 与 $Y$ 为具可数基的局部紧 Hausdorff 空间。设 $\mu\in\mathscr{R}^+_f(X)$，$u:X\to Y$ 为 $(\mathscr{B}(X),\mathscr{B}(Y))$-可测，$\nu_u\in\mathcal{Y}^+(X;Y)$ 为与 $u$ 相伴的 Young 测度。则
\begin{align}
(\nu_u)_x=\delta_{u(x)}
\quad\text{对 }\mu\text{-几乎所有 }x\in X.\label{align:ch05-5.53}
\end{align}
\end{proposition}

\begin{proof}
对任意 $F\in\mathscr{B}(Y)$，由\eqref{align:ch05-5.36}得
\begin{align*}
\nu_u(X\times F)=\int_{X\times Y}\chi_F(y)\,\mathrm{d}\nu_u(x,y)
=\int_X(\nu_u)_x(F)\,\mathrm{d}\mu(x).
\end{align*}
另一方面，由定义
\begin{align*}
\nu_u(X\times F)=\mu(u^{-1}(F))
=\int_X\chi_{u^{-1}(F)}(x)\,\mathrm{d}\mu(x).
\end{align*}
由\eqref{align:ch02-2.29}，对每个 $F\in\mathscr{B}(Y)$，存在 $\mu$-零测集 $N\subseteq X$ 使对所有 $x\in N^c$，
\begin{align*}
(\nu_u)_x(F)=\chi_{u^{-1}(F)}(x)=
\begin{cases}
1 & \text{若 }u(x)\in F,\\
0 & \text{否则}.
\end{cases}
\end{align*}
由于 $X$ 具有可数基，可测空间 $(X,\mathscr{B}(X))$ 可分（见\cref{definition:ch01-可分空间}），故由标准论证，$\mu$-零测集 $N$ 的选取可独立于 $F\in\mathscr{B}(Y)$。结论得证。
\end{proof}

由等式\eqref{align:ch05-5.36}与\eqref{align:ch05-5.53}立即得如下结果。

\begin{corollary}\label{corollary:ch05-Young积分}
设 $X$ 与 $Y$ 为具可数基的局部紧 Hausdorff 空间。设 $\mu\in\mathscr{R}^+_f(X)$，$u:X\to Y$ 为 $(\mathscr{B}(X),\mathscr{B}(Y))$-可测，$\nu_u\in\mathcal{Y}^+(X,\mathscr{B}(X),\mu;Y)$ 为与 $u$ 相伴的 Young 测度。则对每个 $h\in L^1(X\times Y,\mathscr{B}(X\times Y),\nu_u)$，
\begin{align}
\int_{X\times Y}h(x,y)\,\mathrm{d}\nu_u(x,y)
=\int_X h(x,u(x))\,\mathrm{d}\mu(x).\label{align:ch05-5.54}
\end{align}
\end{corollary}

鉴于\cref{proposition:ch05-离散分解}，考虑从 $X$ 到 $\mathscr{R}_f(Y)$ 的弱$^*$可测函数集是有趣的，这些函数对 $\mu$-几乎所有 $x\in X$ 取 $Y$ 上概率测度集 $P(Y)\subseteq\mathscr{R}^+_f(Y)$ 中的值。我们记 $P(X;Y)\equiv P(X,\mathscr{B}(X),\mu;Y)$ 为该集，称为 $Y$ 上的\textbf{参数化概率}集。易见 $P(X;Y)$ 是空间 $L^\infty_{w^*}(X;\mathscr{R}_f(Y))$ 的单位球
\begin{align}
U:=\left\{\nu\in L^\infty_{w^*}(X;\mathscr{R}_f(Y))\ \middle|\
\operatorname*{ess\,sup}_{x\in X}\|\nu_x\|_{\mathscr{R}_f(Y)}\leqslant1\right\}\label{align:ch05-5.55}
\end{align}
的闭凸子集。

\begin{corollary}\label{corollary:ch05-Young一一对应}
设 $X$ 与 $Y$ 为具可数基的局部紧 Hausdorff 空间，$\mu\in\mathscr{R}^+_f(X)$。则 $\mathcal{Y}^+(X;Y)$ 与 $P(X;Y)$ 之间存在一一对应。
\end{corollary}

\begin{proof}
由\cref{proposition:ch05-离散分解}，每个 $\nu\in\mathcal{Y}^+(X;Y)$ 对应 $Y$ 上的一个参数化概率。反之，设 $\{\nu_x\}_{x\in X}\in P(X;Y)$，置
\begin{align*}
\nu:\mathscr{B}(X\times Y)\to[0,\infty),\quad
\nu(G):=\int_X\nu_x(G_x)\,\mathrm{d}\mu(x)
\quad(G\in\mathscr{B}(X\times Y))
\end{align*}
（注意由\cref{remark:ch05-参数测度等价}该定义良定）。由\cref{proposition:ch01-Radon投影}，$\nu\in\mathscr{R}^+_f(X\times Y)$，取 $G=E\times Y$（$E\in\mathscr{B}(X)$）得
\begin{align*}
\nu(E\times Y)=\int_E\nu_x(Y)\,\mathrm{d}\mu(x)=\mu(E).
\end{align*}
故 $\nu\in\mathcal{Y}^+(X;Y)$（见\cref{definition:ch02-Young测度}），结论得证。
\end{proof}

由\cref{corollary:ch05-Young一一对应}与前述注记，自然赋予 $\mathcal{Y}^+(X;Y)$ 以 $L^\infty_{w^*}(X;\mathscr{R}_f(Y))$ 上弱$^*$收敛的拓扑（见\cref{subsection:ch05-Young弱星}）。$\mathcal{Y}^+(X;Y)$ 上另一种与窄收敛相关的拓扑将在\cref{subsection:ch05-Young窄}中引入。为此，我们需要如下定义。

\begin{definition}\label{definition:ch05-Caratheodory}
设 $(X,\mathcal{A},\mu)$ 为有限测度空间，$Y$ 为拓扑空间。
\begin{enumerate}[(i)]
\item\label{item:ch05-car-i} \textbf{Carathéodory 函数}是指任意 $h:X\times Y\to\mathbb{R}$ 满足
\begin{align*}
\begin{cases}
\text{(a) } h(x,\cdot)\in C(Y) & \text{对 }\mu\text{-几乎所有 }x\in X;\\
\text{(b) } h(\cdot,y)\text{ 为 }\mathcal{A}\text{-可测} & \text{对所有 }y\in Y.
\end{cases}
\end{align*}
\item\label{item:ch05-car-ii} \textbf{有界 Carathéodory 被积函数}是指任意 Carathéodory 函数满足
\begin{align*}
\begin{cases}
\text{(a) } h(x,\cdot)\in C_b(Y) & \text{对 }\mu\text{-几乎所有 }x\in X;\\
\text{(b) 映射 }x\mapsto\|h(x,\cdot)\|_\infty\text{ 为 }\mathcal{A}\text{-可测},
\text{ 且 }\int_X\|h(x,\cdot)\|_\infty\,\mathrm{d}\mu(x)<\infty.
\end{cases}
\end{align*}
\end{enumerate}
$C_0$-Carathéodory 被积函数由在 (a) 中将 $C_b(Y)$ 替换为 $C_0(Y)$ 定义。

两个 Carathéodory 函数 $h$ 与 $h'$ 属于同一等价类，若 $h(x,\cdot)=h'(x,\cdot)$ 对 $\mu$-几乎所有 $x\in X$ 成立。
\end{definition}

Carathéodory 函数（等价类）所成的向量空间记作 $\mathcal{C}(X\times Y)$，有界 Carathéodory 被积函数与 $C_0$-Carathéodory 被积函数所成的空间分别记作 $\mathcal{C}_b(X\times Y)$ 与 $\mathcal{C}_0(X\times Y)$。

立即可见，映射 $\mathcal{C}_b(X\times Y)\to[0,\infty)$，
\begin{align}
h\mapsto\|h\|_{\mathcal{C}_b(X\times Y)}
:=\int_X\|h(x,\cdot)\|_\infty\,\mathrm{d}\mu(x),\label{align:ch05-5.56}
\end{align}
是一个范数。于是我们有如下结果。

\begin{proposition}\label{proposition:ch05-Caratheodory同构}
设 $(X,\mathcal{A},\mu)$ 为有限测度空间，$Y$ 为具可数基的局部紧 Hausdorff 空间。则映射
\begin{align}
T:\mathcal{C}_0(X\times Y)\to L^1(X;C_0(Y)),\qquad
(Th)(x):=h(x,\cdot)
\quad\text{对 }\mu\text{-几乎所有 }x\in X,\label{align:ch05-5.57}
\end{align}
是 $(\mathcal{C}_0(X\times Y),\|\cdot\|_{\mathcal{C}_0(X\times Y)})$ 到 $L^1(X;C_0(Y))$ 的等距同构。特别地，$\mathcal{C}_0(X\times Y)$ 赋以范数 $\|\cdot\|_{\mathcal{C}_b(X\times Y)}$ 为 Banach 空间。
\end{proposition}

\begin{proof}
我们先证算子 $T$ 良定。对任意 $h\in\mathcal{C}_0(X\times Y)$ 与 $g\in C_0(Y)$，函数 $(x,y)\mapsto h(x,y)-g(y)$ 属于 $\mathcal{C}_0(X\times Y)$，故对每个 $\alpha\in\mathbb{R}$，集合 $\{x\in X\mid\|h(x,\cdot)-g\|_\infty<\alpha\}$ 属于 $\mathcal{A}$（见\cref{definition:ch05-Caratheodory}(ii)）。另一方面，$C_0(Y)$ 可分（见附录 A,\cref{subsection:appA-可分}），故 Borel $\sigma$-代数 $\mathscr{B}\equiv\mathscr{B}(C_0(Y))$ 由可数开球基 $B(g_k,\alpha):=\{g\in C_0(Y)\mid\|g-g_k\|_\infty<\alpha\}$（$k\in\mathbb{N}$）生成。由此映射 $x\mapsto h(x,\cdot)=(Th)(x)$ 为 $(\mathcal{A},\mathscr{B})$-可测，因而由\cref{proposition:ch04-可测性}为 $\mu$-可测，因为 $C_0(Y)$ 可分。此外，
\begin{align}
\|h\|_{\mathcal{C}_b(X\times Y)}
=\int_X\|h(x,\cdot)\|_\infty\,\mathrm{d}\mu(x)
=\int_X\|(Th)(x)\|_\infty\,\mathrm{d}\mu(x)
=\|Th\|_{L^1(X;C_0(Y))}<\infty.\label{align:ch05-5.58}
\end{align}
故 $Th\in L^1(X;C_0(Y))$，断言得证。

由于 $T$ 线性且等距（见\eqref{align:ch05-5.58}），它单射。为证满射，固定任意 $\hat{h}\in L^1(X;C_0(Y))$。则对 $\mu$-几乎所有 $x\in X$，$\hat{h}(x)\in C_0(Y)$，且映射 $x\mapsto\|\hat{h}(x)\|_\infty$ 属于 $L^1(X)$。此外，对任意 $y\in Y$，映射 $x\mapsto\hat{h}(x)(y)$ 为 $\mathcal{A}$-可测，因为映射 $x\mapsto\|\hat{h}(x)\|_\infty$ 为 $\mathcal{A}$-可测，且映射 $C_0(Y)\to\mathbb{R}$，$g\mapsto g(y)$（$g\in C_0(Y)$，$y\in Y$）连续。因此函数 $(x,y)\mapsto\hat{h}(x)(y)$ 属于 $\mathcal{C}_0(X\times Y)$，故 $T$ 满射。结论得证。
\end{proof}

由\cref{proposition:ch05-Caratheodory同构}、\cref{theorem:ch04-可分}与\cref{proposition:ch04-对偶}得如下结果。

\begin{corollary}\label{corollary:ch05-Caratheodory对偶}
设 $(X,\mathcal{A},\mu)$ 为有限测度空间，$Y$ 为具可数基的局部紧 Hausdorff 空间。则：
\begin{enumerate}[(i)]
\item\label{item:ch05-cd-i} 空间 $\mathcal{C}_0(X\times Y)$ 可分；
\item\label{item:ch05-cd-ii} 我们有
\begin{align*}
(\mathcal{C}_0(X\times Y))^*=L^\infty_{w^*}(X;\mathscr{R}_f(Y)).
\end{align*}
\end{enumerate}
\end{corollary}

\begin{remark}\label{remark:ch05-紧空间同构}
设 $(X,\mathcal{A},\mu)$ 为有限测度空间，$Y$ 为紧空间。则 $C_b(Y)$ 可分，且\cref{proposition:ch05-Caratheodory同构}的证明表明映射
\begin{align}
T:\mathcal{C}_b(X\times Y)\to L^1(X;C_b(Y)),\qquad
(Th)(x):=h(x,\cdot)
\quad\text{对 }\mu\text{-几乎所有 }x\in X,\label{align:ch05-5.59}
\end{align}
是 $(\mathcal{C}_b(X\times Y),\|\cdot\|_{\mathcal{C}_b(X\times Y)})$ 到 $L^1(X;C_b(Y))$ 的等距同构。
\end{remark}

\subsection{弱$^*$收敛}\label{subsection:ch05-Young弱星}

如已观察到的，由\cref{corollary:ch05-Young一一对应}与前述注记，自然赋予 $\mathcal{Y}^+(X;Y)$ 以 $L^\infty_{w^*}(X;\mathscr{R}_f(Y))$ 上弱$^*$收敛的拓扑。于是由\cref{corollary:ch05-Caratheodory对偶}(ii)得如下定义（回忆由定义 $\mathcal{Y}^+(X;Y)\subseteq\mathscr{R}^+_f(X\times Y)$）。

\begin{definition}\label{definition:ch05-Young弱星}
设 $X$ 与 $Y$ 为具可数基的局部紧 Hausdorff 空间。序列 $\{\nu_k\}\subseteq\mathcal{Y}^+(X;Y)$ \textbf{弱$^*$收敛}于 $\nu\in L^\infty_{w^*}(X;\mathscr{R}_f(Y))$（记作 $\nu_k\stackrel{*}{\rightharpoonup}\nu$），若
\begin{align*}
\lim_{k\to\infty}\int_{X\times Y}h\,\mathrm{d}\nu_k
=\int_{X\times Y}h\,\mathrm{d}\nu
\quad\text{对所有 }h\in\mathcal{C}_0(X\times Y).
\end{align*}
$\mathcal{Y}^+(X;Y)$ 上的弱$^*$拓扑，记作 $\mathcal{T}_{w^*}$，是使所有映射 $\nu\mapsto\int_{X\times Y}h\,\mathrm{d}\nu$（$h\in\mathcal{C}_0(X\times Y)$）连续的最弱拓扑。

$\mathcal{T}_{w^*}$ 的一组基为子集族
\begin{align}
B_{h_1,\dots,h_n;\epsilon}(\nu_0)
:=\left\{\nu\in\mathcal{Y}^+(X;Y)\ \middle|\
\left|\int_{X\times Y}h_j\,\mathrm{d}\nu-\int_{X\times Y}h_j\,\mathrm{d}\nu_0\right|<\epsilon
\ \forall j=1,\dots,n\right\},\label{align:ch05-5.60}
\end{align}
其中 $h_1,\dots,h_n\in\mathcal{C}_0(X\times Y)$，$\nu_0\in\mathcal{Y}^+(X;Y)$，$n\in\mathbb{N}$，$\epsilon>0$。
\end{definition}

\begin{remark}\label{remark:ch05-弱星损失质量}
由\cref{definition:ch05-Young弱星}与 $L^\infty_{w^*}(X;\mathscr{R}_f(Y))$ 中范数的下半连续性，弱$^*$极限 $\nu$ 属于单位球 $U\subseteq L^\infty_{w^*}(X;\mathscr{R}_f(Y))$（见\eqref{align:ch05-5.55}）。然而，由于可能的“质量损失”，$\nu$ 不必是 $Y$ 上的参数化概率（见\cref{remark:ch05-范数下半连续}）。还注意 $\nu\in L^\infty_{w^*}(X;\mathscr{R}^+_f(Y))$，因为 $\{\nu_k\}\subseteq\mathcal{Y}^+(X;Y)$。
\end{remark}

\begin{theorem}\label{theorem:ch05-Young弱星紧}
设 $X$ 与 $Y$ 为具可数基的局部紧 Hausdorff 空间。则对每个序列 $\{\nu_k\}\subseteq\mathcal{Y}^+(X;Y)$，存在子列 $\{\nu_{k_l}\}\subseteq\{\nu_k\}$ 与 $\nu\in U\subseteq L^\infty_{w^*}(X;\mathscr{R}_f(Y))$ 使 $\nu_{k_l}\stackrel{*}{\rightharpoonup}\nu$。
\end{theorem}

\begin{proof}
由\cref{theorem:ch04-可分}，空间 $L^1(X;C_0(Y))$ 可分，故由\cref{proposition:ch04-对偶}与 Banach 定理，单位球 $U\subseteq L^\infty_{w^*}(X;\mathscr{R}_f(Y))$ 弱$^*$紧。于是由\cref{corollary:ch05-Caratheodory对偶}结论得证。
\end{proof}

\begin{remark}\label{remark:ch05-Young弱星度量}
由于 $L^1(X;C_0(Y))=\mathcal{C}_0(X\times Y)$ 可分，相对弱$^*$拓扑 $\mathcal{T}_{w^*}\cap U$ 可度量化。事实上，由\cref{proposition:appA-弱星可度量化}，它与由范数
\begin{align}
\nu\mapsto|||\nu|||
:=\sum_{k=1}^\infty\frac{1}{2^k}
\frac{\left|\int_{X\times Y}h_k\,\mathrm{d}\nu\right|}
{1+\int_X\|h_k(x,\cdot)\|_\infty\,\mathrm{d}\mu(x)}\label{align:ch05-5.61}
\end{align}
在 $L^\infty_{w^*}(X;\mathscr{R}_f(Y))$ 上诱导的度量拓扑一致，其中 $D\equiv\{h_k\}\subseteq\mathcal{C}_0(X\times Y)$ 为可数稠密子集。
\end{remark}

\subsection{窄收敛与胎紧性}\label{subsection:ch05-Young窄}

\begin{definition}\label{definition:ch05-Young窄}
设 $X$ 与 $Y$ 为具可数基的局部紧 Hausdorff 空间。序列 $\{\nu_k\}\subseteq\mathcal{Y}^+(X;Y)$ \textbf{窄收敛}于 $\nu\in\mathcal{Y}^+(X;Y)$（记作 $\nu_k\stackrel{n}{\rightharpoonup}\nu$），若
\begin{align*}
\lim_{k\to\infty}\int_{X\times Y}h\,\mathrm{d}\nu_k
=\int_{X\times Y}h\,\mathrm{d}\nu
\quad\text{对所有 }h\in\mathcal{C}_b(X\times Y).
\end{align*}
$\mathcal{Y}^+(X;Y)$ 上的窄拓扑，记作 $\mathcal{T}_n$，是使所有映射 $\nu\mapsto\int_{X\times Y}h\,\mathrm{d}\nu$（$h\in\mathcal{C}_b(X\times Y)$）连续的最弱拓扑。

$\mathcal{T}_n$ 的一组基为子集族\eqref{align:ch05-5.60}，只是其中 $h_1,\dots,h_n\in\mathcal{C}_b(X\times Y)$。
\end{definition}

如下半连续性结果将被使用。

\begin{proposition}\label{proposition:ch05-Caratheodory下半连续}
设 $X$ 与 $Y$ 为具可数基的局部紧 Hausdorff 空间，$h:X\times Y\to[0,\infty)$ 为 Carathéodory 函数。则映射 $\mathcal{Y}^+(X;Y)\to[0,\infty]$，$\nu\mapsto\int_{X\times Y}h\,\mathrm{d}\nu$，在窄拓扑中下半连续。
\end{proposition}

\begin{proof}
我们须证明：若在 $\mathcal{Y}^+(X;Y)$ 中 $\nu_k\stackrel{n}{\rightharpoonup}\nu$，则
\begin{align}
\int_{X\times Y}h(x,y)\,\mathrm{d}\nu(x,y)
\leqslant\liminf_{k\to\infty}\int_X\mathrm{d}\mu(x)\int_Y h(x,y)\,\mathrm{d}(\nu_k)_x(y).\label{align:ch05-5.62}
\end{align}
设 $d$ 为空间 $Y$ 上的相容度量（由\cref{proposition:appA-紧生成}，$Y$ 可度量化）。对每个 $n\in\mathbb{N}$，定义 $h_n:X\times Y\to[0,\infty)$ 为
\begin{align*}
h_n(x,y):=\min\left\{\inf_{\zeta\in Y}[h(x,\zeta)+n\,d(\zeta,y)],\ n\right\}
\qquad(x\in X,\ y\in Y).
\end{align*}
易见对每个 $n\in\mathbb{N}$，$h_n$ 为 $\mathscr{B}(X\times Y)$-可测，$h_n\in\mathcal{C}_b(X\times Y)$，$h_n\leqslant h_{n+1}$，且 $h_n\to h$ 在 $X\times Y$ 上逐点收敛。于是由单调收敛定理，
\begin{align*}
\int_{X\times Y}h\,\mathrm{d}\nu
=\lim_{n\to\infty}\int_{X\times Y}h_n\,\mathrm{d}\nu
=\sup_{n\in\mathbb{N}}\int_{X\times Y}h_n\,\mathrm{d}\nu.
\end{align*}
对每个 $n\in\mathbb{N}$，映射 $\nu\mapsto\int_{X\times Y}h_n\,\mathrm{d}\nu$ 由\cref{definition:ch05-Young窄}连续，因为 $h_n\in\mathcal{C}_b(X\times Y)$。故映射 $\nu\mapsto\int_{X\times Y}h\,\mathrm{d}\nu$ 作为连续函数序列的上确界为下半连续。结果得证。
\end{proof}

对 Young 测度，胎紧性定义如下。

\begin{definition}\label{definition:ch05-Young胎紧}
设 $X$ 与 $Y$ 为具可数基的局部紧 Hausdorff 空间。子集 $N\subseteq\mathcal{Y}^+(X;Y)$ 称为\textbf{胎紧的}，若对任意 $\epsilon>0$，存在紧子集 $K\subseteq Y$ 使对所有 $\nu\in N$ 有 $\nu(X\times K^c)<\epsilon$。
\end{definition}

\begin{remark}\label{remark:ch05-Young胎紧投影}
由\cref{definition:ch05-Young胎紧}，子集 $N\subseteq\mathcal{Y}^+(X;Y)$ 胎紧，当且仅当投影到 $Y$ 的集合 $\{\nu_{p_Y}\mid\nu\in N\}\subseteq\mathscr{R}^+_f(Y)$ 在\cref{definition:ch05-胎紧}意义下胎紧。
\end{remark}

如下结果的证明与\cref{proposition:ch05-胎紧拓扑一致}严格类似，故我们略去。

\begin{proposition}\label{proposition:ch05-Young胎紧一致}\label{proposition:ch05-矩估计}
设 $X$ 与 $Y$ 为具可数基的局部紧 Hausdorff 空间，$N\subseteq\mathcal{Y}^+(X;Y)$ 胎紧。则 $N$ 上的窄拓扑与 $N$ 上的弱$^*$拓扑一致，即 $\mathcal{T}_n\cap N=\mathcal{T}_{w^*}\cap N$。
\end{proposition}

作为胎紧性判据，由\cref{remark:ch05-Young胎紧投影}与\cref{proposition:ch05-胎紧判据}得如下结果。

\begin{proposition}\label{proposition:ch05-Young胎紧判据}
设 $X$ 与 $Y$ 为具可数基的局部紧 Hausdorff 空间。则下列命题等价：
\begin{enumerate}[(i)]
\item\label{item:ch05-ytight-i} $N\subseteq\mathcal{Y}^+(X;Y)$ 胎紧；
\item\label{item:ch05-ytight-ii} 存在下紧函数 $\varphi:Y\to[0,\infty)$ 使
\begin{align}
\sup_{\nu\in N}\int_Y\varphi\,\mathrm{d}\nu_{p_Y}<\infty.\label{align:ch05-5.63}
\end{align}
\end{enumerate}
\end{proposition}

对 Young 测度，我们有\cref{theorem:ch05-prokhorov}的如下对应物。

\begin{theorem}\label{theorem:ch05-弱紧Young}\label{proposition:ch05-弱紧YoungL1}
设 $X$ 与 $Y$ 为具可数基的局部紧 Hausdorff 空间。设 $N\subseteq\mathcal{Y}^+(X;Y)$。则下列命题等价：
\begin{enumerate}[(i)]
\item\label{item:ch05-youngcomp-i} $N$ 在窄拓扑中相对序列紧；
\item\label{item:ch05-youngcomp-ii} $N$ 胎紧。
\end{enumerate}
\end{theorem}

\begin{proof}
(i)$\Rightarrow$(ii)。由假设，对任意序列 $\{\nu_k\}\subseteq N$，存在子列 $\{\nu_{k_j}\}\subseteq\{\nu_k\}$ 与 $\nu\in\mathcal{Y}^+(X;Y)$ 使
\begin{align}
\lim_{j\to\infty}\int_{X\times Y}h\,\mathrm{d}\nu_{k_j}
=\int_{X\times Y}h\,\mathrm{d}\nu
\quad\text{对所有 }h\in\mathcal{C}_b(X\times Y).\label{align:ch05-5.64}
\end{align}
取 $h(x,y)=g(y)$，$g\in C_b(Y)$。由于
\begin{align*}
\int_{X\times Y}g(y)\,\mathrm{d}\nu(x,y)=\int_Y g(y)\,\mathrm{d}\nu_{p_Y}(y),\\
\int_{X\times Y}g(y)\,\mathrm{d}\nu_{k_j}(x,y)=\int_Y g(y)\,\mathrm{d}(\nu_{k_j})_{p_Y}(y)
\end{align*}
（见\cref{remark:ch02-Young投影}(i)），由\eqref{align:ch05-5.64}得
\begin{align*}
\lim_{j\to\infty}\int_Y g(y)\,\mathrm{d}(\nu_{k_j})_{p_Y}(y)
=\int_Y g(y)\,\mathrm{d}\nu_{p_Y}(y)
\quad\text{对所有 }g\in C_b(Y).
\end{align*}
由此集合 $\{\nu_{p_Y}\mid\nu\in N\}$ 在 $\mathscr{R}^+_f(Y)$ 的窄拓扑中相对序列紧，故由\cref{theorem:ch05-prokhorov}胎紧。于是由\cref{remark:ch05-Young胎紧投影}断言得证。

(ii)$\Rightarrow$(i)。由\cref{remark:ch02-Young投影}，$N$ 是 $\mathscr{R}^+_f(X\times Y)$ 的有界子集。我们证明它也是 $\mathscr{R}^+_f(X\times Y)$ 的胎紧子集，即对任意 $\epsilon>0$，存在紧子集 $\tilde{K}\subseteq X\times Y$ 使对所有 $\nu\in N$ 有 $\nu(\tilde{K}^c)<\epsilon$（见\cref{definition:ch05-胎紧}）。由假设，对任意 $\epsilon>0$，存在紧子集 $K\subseteq Y$ 使对所有 $\nu\in N$ 有 $\nu_{p_Y}(K^c)=\nu(X\times K^c)<\epsilon/2$。此外，由\cref{proposition:appA-紧生成}，$X$ 为 $\sigma$-紧，故存在紧子集 $K_0\subseteq X$ 使 $\mu(K_0^c)<\epsilon/2$。注意由 Young 测度定义，对每个 $\nu\in N$ 有 $\nu_{p_X}(K_0^c)=\nu(K_0^c\times Y)=\mu(K_0^c)$。置 $\tilde{K}:=K_0\times K\subseteq X\times Y$。显然 $\tilde{K}$ 紧，且 $\tilde{K}^c\subseteq(X\times K^c)\cup(K_0^c\times Y)$。则对所有 $\nu\in N$，
\begin{align*}
\nu(\tilde{K}^c)\leqslant\nu(X\times K^c)+\nu(K_0^c\times Y)
=\nu_{p_Y}(K^c)+\mu(K_0^c)<\epsilon,
\end{align*}
断言得证。

由上述注记与\cref{proposition:ch05-有界胎紧紧}（另见\cref{remark:ch05-polish}），对任意序列 $\{\nu_k\}\subseteq N$，存在子列 $\{\nu_{k_j}\}\subseteq\{\nu_k\}$ 与 $\nu\in\mathscr{R}^+_f(X\times Y)$ 使
\begin{align}
\lim_{j\to\infty}\int_{X\times Y}g\,\mathrm{d}\nu_{k_j}
=\int_{X\times Y}g\,\mathrm{d}\nu
\quad\text{对所有 }g\in C_b(X\times Y).\label{align:ch05-5.65}
\end{align}
易见 $\nu\in\mathcal{Y}^+(X;Y)$。事实上，由于 $\{\nu_{k_j}\}\subseteq N\subseteq\mathcal{Y}^+(X;Y)$，由离散分解定理，等式\eqref{align:ch05-5.65}读作
\begin{align*}
\lim_{j\to\infty}\int_X\mathrm{d}\mu(x)\int_Y g(x,y)\,\mathrm{d}(\nu_{k_j})_x(y)
=\int_X\mathrm{d}\nu_{p_X}(x)\int_Y g(x,y)\,\mathrm{d}\nu_x(y)
\end{align*}
对所有 $g\in\mathcal{C}_b(X\times Y)$。在此等式中取 $g=g(x)$，得
\begin{align*}
\int_X g(x)\,\mathrm{d}\mu(x)=\int_X g(x)\,\mathrm{d}\nu_{p_X}(x)
\quad\text{对所有 }g\in C_b(X).
\end{align*}
由\cref{proposition:ch04-稠密}(ii)（另见\cref{proposition:appA-紧生成}），由此等式显然得 $\nu_{p_X}(E)=\mu(E)$ 对所有 $E\in\mathscr{B}(X)$。断言得证。

特别地，等式\eqref{align:ch05-5.65}对所有 $g\in\mathcal{C}_0(X\times Y)$ 成立。由于空间 $\mathcal{C}_0(X\times Y)$ 在 $L^1(X;C_0(Y))$ 中稠密（见\cref{proposition:ch04-稠密}(ii)与\cref{proposition:ch05-Caratheodory同构}），对任意 $h\in L^1(X;C_0(Y))$，存在序列 $\{g_m\}\subseteq\mathcal{C}_0(X\times Y)$ 使 $\lim_{m\to\infty}\|g_m-h\|_{L^1(X;C_0(Y))}=0$。固定 $\epsilon>0$，取 $m_0\in\mathbb{N}$ 充分大使 $\|g_{m_0}-h\|_{L^1(X;C_0(Y))}<\epsilon/3$。另一方面，由\eqref{align:ch05-5.65}存在 $j_0\in\mathbb{N}$ 使
\begin{align*}
\left|\int_{X\times Y}g_{m_0}\,\mathrm{d}\nu_{k_j}-\int_{X\times Y}g_{m_0}\,\mathrm{d}\nu\right|<\frac{\epsilon}{3}
\quad\text{对所有 }j>j_0.
\end{align*}
则对所有 $j>j_0$，
\begin{align*}
\left|\int_{X\times Y}h\,\mathrm{d}\nu_{k_j}-\int_{X\times Y}h\,\mathrm{d}\nu\right|
\leqslant\left|\int_{X\times Y}(h-g_{m_0})\,\mathrm{d}\nu_{k_j}\right|
+\left|\int_{X\times Y}(h-g_{m_0})\,\mathrm{d}\nu\right|
+\left|\int_{X\times Y}g_{m_0}\,\mathrm{d}\nu_{k_j}-\int_{X\times Y}g_{m_0}\,\mathrm{d}\nu\right|\\
<\left|\int_{X\times Y}(h-g_{m_0})\,\mathrm{d}\nu_{k_j}\right|
+\left|\int_{X\times Y}(h-g_{m_0})\,\mathrm{d}\nu\right|
+\frac{\epsilon}{3}.
\end{align*}
由于 $\{\nu_{k_j}\}\subseteq N\subseteq\mathcal{Y}^+(X;Y)$，由离散分解定理得
\begin{align*}
\left|\int_{X\times Y}(h-g_{m_0})\,\mathrm{d}\nu_{k_j}\right|
=\left|\int_X\mathrm{d}\mu(x)\int_Y(h-g_{m_0})(x,y)\,\mathrm{d}(\nu_{k_j})_x(y)\right|
\leqslant\|g_{m_0}-h\|_{L^1(X;C_0(Y))}<\frac{\epsilon}{3}.
\end{align*}
由于 $\nu\in\mathcal{Y}^+(X;Y)$，我们类似地得
\begin{align*}
\left|\int_{X\times Y}(h-g_{m_0})\,\mathrm{d}\nu\right|
=\left|\int_X\mathrm{d}\mu(x)\int_Y(h-g_{m_0})(x,y)\,\mathrm{d}\nu_x(y)\right|
\leqslant\|g_{m_0}-h\|_{L^1(X;C_0(Y))}<\frac{\epsilon}{3}.
\end{align*}

鉴于上述不等式，对任意 $\epsilon>0$，存在 $j_0\in\mathbb{N}$ 使对所有 $j>j_0$，
\begin{align*}
\left|\int_{X\times Y}h\,\mathrm{d}\nu_{k_j}-\int_{X\times Y}h\,\mathrm{d}\nu\right|<\epsilon
\quad\text{对所有 }h\in\mathcal{C}_0(X\times Y).
\end{align*}
综上，我们证明了：对任意序列 $\{\nu_k\}\subseteq N$，存在子列 $\{\nu_{k_j}\}\subseteq\{\nu_k\}$ 与 $\nu\in\mathcal{Y}^+(X;Y)$ 使 $\nu_{k_j}\stackrel{*}{\rightharpoonup}\nu$。由于 $N$ 胎紧，由\cref{proposition:ch05-Young胎紧一致}结论得证。
\end{proof}

如下结果类似地证明。

\begin{proposition}\label{proposition:ch05-逐点窄收敛}
设 $X$ 与 $Y$ 为具可数基的局部紧 Hausdorff 空间。设 $\{\nu_k\}\subseteq\mathcal{Y}^+(X;Y)$，$\nu\in\mathcal{Y}^+(X;Y)$ 满足对 $\mu$-几乎所有 $x\in X$，$(\nu_k)_x\stackrel{n}{\rightharpoonup}\nu_x$ 在 $\mathscr{R}^+_f(Y)$ 中。则 $\nu_k\stackrel{n}{\rightharpoonup}\nu$ 在 $\mathcal{Y}^+(X;Y)$ 中。
\end{proposition}

\begin{proof}
设 $h\in\mathcal{C}_b(X\times Y)$。由\cref{definition:ch05-Caratheodory}(ii)，对 $\mu$-几乎所有 $x\in X$ 有 $h(x,\cdot)\in C_b(Y)$。则对 $\mu$-几乎所有 $x\in X$，
\begin{align*}
\lim_{k\to\infty}\int_Y h(x,y)\,\mathrm{d}(\nu_k)_x(y)
=\int_Y h(x,y)\,\mathrm{d}\nu_x(y)
\end{align*}
（见\cref{definition:ch05-窄收敛}）。另一方面，对所有 $k\in\mathbb{N}$，
\begin{align*}
\left|\int_Y h(x,y)\,\mathrm{d}(\nu_k)_x(y)\right|
\leqslant\|h(x,\cdot)\|_\infty(\nu_k)_x(Y)=\|h(x,\cdot)\|_\infty,
\end{align*}
且由\cref{definition:ch05-Caratheodory}(ii)，映射 $x\mapsto\|h(x,\cdot)\|_\infty$ 属于 $L^1(X)$。于是由控制收敛定理，
\begin{align*}
\lim_{k\to\infty}\int_X\mathrm{d}\mu(x)\int_Y h(x,y)\,\mathrm{d}(\nu_k)_x(y)
=\int_X\mathrm{d}\mu(x)\int_Y h(x,y)\,\mathrm{d}\nu_x(y),
\end{align*}
结论得证。
\end{proof}

\section{与函数相伴的 Young 测度序列}\label{section:ch05-收敛}

设 $U\subseteq\mathbb{R}^N$ 为有界开集，$\mu:=\lambda_N|_{\mathscr{B}(U)}$，$\{u_j\}\subseteq L^1(U;\mathbb{R}^d)$（我们回忆 $L^1(U;\mathbb{R}^d)\equiv L^1(U,\mathscr{B}(U),\lambda_N|_{\mathscr{B}(U)};\mathbb{R}^d)$）。在\cref{definition:ch02-Young测度}中，我们定义 $\nu_{u_j}\in\mathcal{Y}^+(U;\mathbb{R}^d)$ 为满足
\begin{align*}
\nu_{u_j}(E\times F):=\mu\big(E\cap u_j^{-1}(F)\big)
\quad\text{对所有 }E\in\mathscr{B}(U),\ F\in\mathscr{B}(\mathbb{R}^d)
\end{align*}
的 Young 测度。换言之，$\nu_{u_j}$ 是与 $u_j$ 相伴的参数化概率测度（见\cref{corollary:ch05-Young一一对应}）。由\cref{proposition:ch05-Dirac测度}，
\begin{align}
(\nu_{u_j})_x=\delta_{u_j(x)}
\quad\text{对 }\mu\text{-几乎所有 }x\in U\ (j\in\mathbb{N}).\label{align:ch05-5.66}
\end{align}

本节中我们要确定的序列 $\{u_j\}\subseteq L^1(U;\mathbb{R}^d)$ 的性质，可由相应 Young 测度序列 $\{\nu_{u_j}\}$ 的弱$^*$或窄收敛导出。这两种收敛在\cref{subsection:ch05-Young弱星}与\cref{subsection:ch05-Young窄}中已为一般 Young 测度引入。我们首先将弱$^*$收敛联系到 $u_j$ 的几乎处处收敛（Biting 引理将在\cref{subsection:ch05-Biting}中处理其窄对应）。

\subsection{$\{\nu_{u_j}\}$ 的弱$^*$收敛}

作为热身，我们证明如下结果。

\begin{proposition}\label{proposition:ch05-几乎处处收敛}
设 $\{u_j\}\subseteq L^1(U;\mathbb{R}^d)$。则下列命题等价：
\begin{enumerate}[(i)]
\item\label{item:ch05-ae-i} 对 $\mu$-几乎所有 $x\in U$，$u_j(x)\to u(x)$ 对某个 $u:U\to\mathbb{R}^d$；
\item\label{item:ch05-ae-ii} 存在 Young 测度 $\nu\in\mathcal{Y}^+(U;\mathbb{R}^d)$ 使 $\nu_{u_j}\stackrel{n}{\rightharpoonup}\nu$，且对 $\mu$-几乎所有 $x\in U$ 有 $\nu_x=\delta_{u(x)}$。
\end{enumerate}
\end{proposition}

\begin{proof}
(i)$\Rightarrow$(ii)。设 $h\in\mathcal{C}_b(U\times\mathbb{R}^d)$。则对 $\mu$-几乎所有 $x\in U$，
\begin{align*}
h(x,u_j(x))\to h(x,u(x))
\end{align*}
（由\cref{definition:ch05-Caratheodory}(ii)，$h(x,\cdot)\in C_b(\mathbb{R}^d)$）。此外，
\begin{align*}
|h(x,u_j(x))|\leqslant\|h(x,\cdot)\|_\infty,
\end{align*}
且由\cref{definition:ch05-Caratheodory}(ii)，$x\mapsto\|h(x,\cdot)\|_\infty$ 属于 $L^1(U)$。故由控制收敛定理、\eqref{align:ch05-5.66}与\cref{corollary:ch05-Young积分}，
\begin{align*}
\lim_{j\to\infty}\int_{U\times\mathbb{R}^d}h\,\mathrm{d}\nu_{u_j}
=\lim_{j\to\infty}\int_U h(x,u_j(x))\,\mathrm{d}\mu(x)
=\int_U h(x,u(x))\,\mathrm{d}\mu(x).
\end{align*}
由\eqref{align:ch05-5.66}（其中 $u_j$ 换为 $u$）、\cref{corollary:ch05-Young积分}与\eqref{align:ch05-5.53}，
\begin{align*}
\int_U h(x,u(x))\,\mathrm{d}\mu(x)
=\int_{U\times\mathbb{R}^d}h\,\mathrm{d}\nu_u
=\int_U\mathrm{d}\mu(x)\int_{\mathbb{R}^d}h(x,y)\,\mathrm{d}(\nu_u)_x(y)
=\int_U\mathrm{d}\mu(x)\int_{\mathbb{R}^d}h(x,y)\,\mathrm{d}\delta_{u(x)}(y).
\end{align*}
于是对任意 $h\in\mathcal{C}_b(U\times\mathbb{R}^d)$，
\begin{align*}
\lim_{j\to\infty}\int_{U\times\mathbb{R}^d}h\,\mathrm{d}\nu_{u_j}
=\int_U\mathrm{d}\mu(x)\int_{\mathbb{R}^d}h(x,y)\,\mathrm{d}\delta_{u(x)}(y)
=\int_{U\times\mathbb{R}^d}h\,\mathrm{d}\nu
\end{align*}
其中 $\nu$ 为满足 $\nu_x=\delta_{u(x)}$（对 $\mu$-几乎所有 $x\in U$）的 Young 测度。由\cref{definition:ch05-Young窄}，$\nu_{u_j}\stackrel{n}{\rightharpoonup}\nu$，断言得证。

(ii)$\Rightarrow$(i)。由\cref{definition:ch05-Young窄}，对任意 $h\in\mathcal{C}_b(U\times\mathbb{R}^d)$，
\begin{align*}
\lim_{j\to\infty}\int_{U\times\mathbb{R}^d}h\,\mathrm{d}\nu_{u_j}
=\int_{U\times\mathbb{R}^d}h\,\mathrm{d}\nu.
\end{align*}
特别地，取 $h(x,y)=g(\zeta)\langle y-x,\zeta\rangle$（$g\in C_c(U)$，$\zeta\in\mathbb{R}^d$），由\eqref{align:ch05-5.66}与\cref{corollary:ch05-Young积分}，对任意 $j\in\mathbb{N}$，
\begin{align*}
\int_{U\times\mathbb{R}^d}h\,\mathrm{d}\nu_{u_j}
=\int_U g(x)\langle u_j(x)-x,\zeta\rangle\,\mathrm{d}\mu(x).
\end{align*}
由于 $\nu_x=\delta_{u(x)}$ 对 $\mu$-几乎所有 $x\in U$，对右端有
\begin{align*}
\int_{U\times\mathbb{R}^d}h\,\mathrm{d}\nu
=\int_U\mathrm{d}\mu(x)\int_{\mathbb{R}^d}h(x,y)\,\mathrm{d}\delta_{u(x)}(y)
=\int_U g(x)\langle u(x)-x,\zeta\rangle\,\mathrm{d}\mu(x).
\end{align*}
于是对任意 $g\in C_c(U)$ 与 $\zeta\in\mathbb{R}^d$，
\begin{align}
\lim_{j\to\infty}\int_U g(x)\langle u_j(x),\zeta\rangle\,\mathrm{d}\mu(x)
=\int_U g(x)\langle u(x),\zeta\rangle\,\mathrm{d}\mu(x).\label{align:ch05-5.67}
\end{align}
由 $g$ 与 $\zeta$ 的任意性，由\eqref{align:ch05-5.67}得 $u_j\stackrel{*}{\rightharpoonup}u$ 在 $(L^\infty(U))^*$ 中（用 $\mathbb{R}^d$-向量记号）。另一方面，由\eqref{align:ch05-5.66}，对任意 $j\in\mathbb{N}$，
\begin{align*}
\int_{U\times\mathbb{R}^d}|y|\,\mathrm{d}\nu_{u_j}
=\int_U|u_j(x)|\,\mathrm{d}\mu(x)
=\|u_j\|_{L^1(U;\mathbb{R}^d)}.
\end{align*}
由于对每个 $j\in\mathbb{N}$ 有 $\nu_{u_j}\in\mathcal{Y}^+(U;\mathbb{R}^d)$，映射 $\nu\mapsto\int_{U\times\mathbb{R}^d}|y|\,\mathrm{d}\nu$ 由\cref{proposition:ch05-Caratheodory下半连续}为下半连续。故
\begin{align*}
\int_U|u(x)|\,\mathrm{d}\mu(x)
=\int_{U\times\mathbb{R}^d}|y|\,\mathrm{d}\nu
\leqslant\liminf_{j\to\infty}\int_{U\times\mathbb{R}^d}|y|\,\mathrm{d}\nu_{u_j}
=\liminf_{j\to\infty}\|u_j\|_{L^1(U;\mathbb{R}^d)}<\infty.
\end{align*}
由此 $u\in L^1(U;\mathbb{R}^d)$，故 $u_j\stackrel{*}{\rightharpoonup}u$ 在 $(L^\infty(U;\mathbb{R}^d))^*$ 中。于是由\cref{theorem:ch02-dunfordpettis}(ii)，$\{u_j\}\subseteq L^1(U;\mathbb{R}^d)$ 有界且一致可积。

综上，在假设 (i) 下，$\{u_j\}$ 弱*收敛于 $u\in L^1(U;\mathbb{R}^d)$ 且一致可积。由于 $u_j\stackrel{*}{\rightharpoonup}u$，由此标准论证（例如见\cite[Theorem 3.10 and Corollary 3.11]{ref14}）得：存在子列 $\{u_{j_k}\}\subseteq\{u_j\}$ 使 $u_{j_k}(x)\to u(x)$ 对 $\mu$-几乎所有 $x\in U$。

于是对任意序列 $\{u_j\}$ 满足 (ii)，存在子列在 $\mu$-几乎处处收敛于 $u$。假设 (i) 不成立，则存在 $\{u_j\}$ 满足 (ii) 但无任何子列 $\mu$-几乎处处收敛于 $u$。由此可构造子列 $\{u_{j_k}\}$（仍满足 (ii)）使对 $\mu$-几乎所有 $x\in U$，$\{u_{j_k}(x)\}\subseteq\mathbb{R}^d$ 不收敛于 $u(x)$。但这与上段应用于 $\{u_{j_k}\}$ 矛盾。故 (i) 成立，证明完成。
\end{proof}

\cref{proposition:ch05-几乎处处收敛}(ii)的如下推论有用。

\begin{remark}\label{remark:ch05-收敛注}
设 $\{u_j\}\subseteq L^1(U;\mathbb{R}^d)$。
\begin{enumerate}[(i)]
\item 若 $\{u_j\}$ 一致有界，且 $u_j(x)\to u(x)$ 对 $\mu$-几乎所有 $x\in U$，则存在 Young 测度 $\nu\in\mathcal{Y}^+(U;\mathbb{R}^d)$ 使 $\nu_{u_j}\stackrel{n}{\rightharpoonup}\nu$，且对 $\mu$-几乎所有 $x\in U$ 有 $\nu_x=\delta_{u(x)}$（见\cref{proposition:ch05-几乎处处收敛}）；
\item 若 $\{u_j\}$ 有界且一致可积，则应用 (i) 于 $h(y)=y$ 表明：对某个子列，$u_{j_k}\rightharpoonup u_*$ 在 $L^1(U;\mathbb{R}^d)$ 中，其中 $u_*(x):=\int_{\mathbb{R}^d}y\,\mathrm{d}\nu_x(y)$。元素 $u_*$ 称为 Young 测度 $\nu$ 的\textbf{重心}。
\end{enumerate}
\end{remark}

结合\cref{theorem:ch05-弱星紧}，\cref{proposition:ch05-几乎处处收敛}给出如下著名结果（见\cite[Theorem 2.1]{ref08}）。

\begin{theorem}[Ball]\label{theorem:ch05-ball}
设 $\{u_j\}\subseteq L^1(U;\mathbb{R}^d)$ 有界。则存在子列 $\{u_{j_k}\}\subseteq\{u_j\}$ 与 Young 测度 $\nu\in\mathcal{Y}^+(U;\mathbb{R}^d)$ 使 $\nu_{u_{j_k}}\stackrel{*}{\rightharpoonup}\nu$。
\end{theorem}

\begin{remark}\label{remark:ch05-ball注}
\cref{theorem:ch05-ball}的假设可如下放松。设 $\{u_j\}\subseteq L^1(U;\mathbb{R}^d)$ 使 $\{F(u_j)\}$ 在 $L^1(U)$ 中有界，其中 $F:\mathbb{R}^d\to[0,\infty]$ 为可测下紧函数（见\cref{definition:ch05-下紧}）。则对某个 $\nu\in\mathcal{Y}^+(U;\mathbb{R}^d)$，$\nu_{u_{j_k}}\stackrel{*}{\rightharpoonup}\nu$。事实上，由\cref{proposition:ch05-弱星紧}存在子列 $\{u_{j_k}\}\subseteq\{u_j\}$ 与 $\nu\in U\subseteq L^\infty_{w^*}(U;\mathscr{R}_f(\mathbb{R}^d))$ 使 $\nu_{u_{j_k}}\stackrel{*}{\rightharpoonup}\nu$。此外，由\cref{proposition:ch05-Caratheodory下半连续}、$F$ 的下紧性与假设，
\begin{align*}
\int_{U\times\mathbb{R}^d}F(y)\,\mathrm{d}\nu
\leqslant\liminf_{k\to\infty}\int_{U\times\mathbb{R}^d}F(y)\,\mathrm{d}\nu_{u_{j_k}}
=\liminf_{k\to\infty}\int_U F(u_{j_k}(x))\,\mathrm{d}\mu(x)<\infty
\end{align*}
（见\cref{corollary:ch05-Young积分}）。由此 $\nu_{p_{\mathbb{R}^d}}(\{\infty\})=0$，故 $\nu\in\mathcal{Y}^+(U;\mathbb{R}^d)$（因为 $\nu\in L^\infty_{w^*}(U;\mathscr{R}^+_f(\mathbb{R}^d))$）。
\end{remark}

\subsection{$\{\nu_{u_j}\}$ 的窄收敛}

在\cref{subsection:ch05-Young窄}中已观察到，窄拓扑在紧子集上与弱$^*$拓扑一致，因为 $\mathcal{Y}^+(U;\mathbb{R}^d)$ 的紧子集在弱$^*$拓扑中闭（见\cref{proposition:ch05-Young胎紧一致}）。由此与\cref{theorem:ch05-ball}相同的论证给出如下结果。

\begin{theorem}\label{theorem:ch05-函数窄收敛}\label{proposition:ch05-5.4.1}\label{proposition:ch05-弱收敛}
设 $\{u_j\}\subseteq L^1(U;\mathbb{R}^d)$ 使 $\{F(u_j)\}$ 在 $L^1(U)$ 中有界，其中 $F:\mathbb{R}^d\to[0,\infty]$ 为下紧函数。则存在子列 $\{u_{j_k}\}\subseteq\{u_j\}$ 与 $\nu\in\mathcal{Y}^+(U;\mathbb{R}^d)$ 使 $\nu_{u_{j_k}}\stackrel{n}{\rightharpoonup}\nu$。
\end{theorem}

\begin{proof}
由\cref{remark:ch05-ball注}存在子列 $\{u_{j_k}\}\subseteq\{u_j\}$ 与 $\nu\in\mathcal{Y}^+(U;\mathbb{R}^d)$ 使 $\nu_{u_{j_k}}\stackrel{*}{\rightharpoonup}\nu$。此外，对任意紧子集 $K\subseteq\mathbb{R}^d$，我们有
\begin{align*}
\sup_{k\in\mathbb{N}}\nu_{u_{j_k}}(U\times K^c)
=\sup_{k\in\mathbb{N}}\int_U\chi_{K^c}(u_{j_k}(x))\,\mathrm{d}\mu(x)
=\sup_{k\in\mathbb{N}}\mu(\{u_{j_k}\in K^c\})
\leqslant\sup_{k\in\mathbb{N}}\frac{\|F(u_{j_k})\|_{L^1(U)}}{\inf_{K^c}F}<\infty
\end{align*}
（回忆 $F$ 下紧，故对某个 $R>0$ 有 $\inf_{K^c}F>R$）。由此由\cref{proposition:ch05-Young胎紧判据}，集合 $\{\nu_{u_{j_k}}\mid k\in\mathbb{N}\}\subseteq\mathcal{Y}^+(U;\mathbb{R}^d)$ 胎紧。于是由\cref{proposition:ch05-Young胎紧一致}得 $\nu_{u_{j_k}}\stackrel{n}{\rightharpoonup}\nu$。结论得证。
\end{proof}

\begin{proposition}\label{proposition:ch05-有界紧}\label{proposition:ch05-弱星紧}
设 $\{u_j\}\subseteq L^1(U;\mathbb{R}^d)$ 使 $\{F(u_j)\}$ 在 $L^1(U)$ 中有界，$F:\mathbb{R}^d\to[0,\infty]$ 下紧。则集族 $\{\nu_{u_j}\mid j\in\mathbb{N}\}\subseteq\mathcal{Y}^+(U;\mathbb{R}^d)$ 胎紧。
\end{proposition}

\begin{remark}\label{remark:ch05-周期}
设 $u\in L^1_\mathrm{loc}(\mathbb{R}^d;\mathbb{R}^d)$ 为 $\mathbb{Z}^d$-周期函数。对每个 $j\in\mathbb{N}$，置 $u_j(x):=u(jx)$（$x\in U$）。则 $\{u_j\}$ 在 $L^1(U;\mathbb{R}^d)$ 中有界。由\cref{theorem:ch05-函数窄收敛}，存在子列 $\{u_{j_k}\}\subseteq\{u_j\}$ 与 $\nu\in\mathcal{Y}^+(U;\mathbb{R}^d)$ 使 $\nu_{u_{j_k}}\stackrel{n}{\rightharpoonup}\nu$。此外，对任意 $h\in\mathcal{C}_b(U\times\mathbb{R}^d)$ 与 $\mathbb{Z}^d$-胞腔 $Q$，
\begin{align*}
\int_U h(x,u_{j_k}(x))\,\mathrm{d}\mu(x)
=\int_U h(x,u(j_kx))\,\mathrm{d}\mu(x)
=\sum_{z\in\mathbb{Z}^d}\int_{j_k^{-1}Q+z/j_k}h(x,u(j_kx))\,\mathrm{d}\mu(x)\\
=\sum_{z\in\mathbb{Z}^d}\int_{Q+z}h(j_k^{-1}(y+z),u(y))\,\mathrm{d}\mu(y)\to\int_U\mathrm{d}\mu(x)\int_Q h(x,u(y))\,\mathrm{d}\mu(y).
\end{align*}
故在此情形，对 $\mu$-几乎所有 $x\in U$ 有 $\nu_x=\nu_u\in\mathcal{Y}^+(Q;\mathbb{R}^d)$。
\end{remark}

如下例子展示与无界序列 $\{u_j\}\subseteq L^1(U;\mathbb{R}^d)$ 相伴的 Young 测度（见\cite[Example 1.1]{ref92}）。

\begin{proposition}\label{proposition:ch05-非相伴}
存在序列 $\{u_j\}\subseteq L^1(0,1;\mathbb{R})$ 与 Young 测度 $\nu\in\mathcal{Y}^+(0,1;\mathbb{R})$ 使 $u_j\rightharpoonup u$ 在 $L^1(0,1)$ 中且 $\nu_{u_j}\stackrel{n}{\rightharpoonup}\nu$，但 $\nu$ 不与任何 $v\in L^1(0,1;\mathbb{R})$ 相伴。
\end{proposition}

\begin{proof}
对每个 $j\in\mathbb{N}$，置
\begin{align*}
u_j(x):=\sin(2\pi jx)\qquad(x\in(0,1)).
\end{align*}
由于 $\{u_j\}$ 在 $L^1(0,1;\mathbb{R})$ 中弱收敛于 $u=0$（例如由 Riemann–Lebesgue 引理）。此外，$\{u_j\}$ 在 $L^\infty(0,1;\mathbb{R})$ 中有界，故由\cref{remark:ch05-收敛注}(i)–(ii)，存在子列 $\{u_{j_k}\}\subseteq\{u_j\}$ 与 Young 测度 $\nu\in\mathcal{Y}^+(0,1;\mathbb{R})$ 使 $\nu_{u_{j_k}}\stackrel{n}{\rightharpoonup}\nu$，且
\begin{align*}
u_*(x):=\int_{\mathbb{R}}y\,\mathrm{d}\nu_x(y)=0
\quad\text{对 }\mu\text{-几乎所有 }x\in(0,1).
\end{align*}
设 $h\in\mathcal{C}_b((0,1)\times\mathbb{R})$，则
\begin{align*}
\int_0^1 h(x,u_j(x))\,\mathrm{d}x=\int_0^1 h(x,\sin(2\pi jx))\,\mathrm{d}x\to\int_0^1\mathrm{d}x\int_0^1 h(x,\sin(2\pi y))\,\mathrm{d}y
\end{align*}
（用\cite[Lemma 2.37]{ref45}的标准论证）。于是对 $\mu$-几乎所有 $x\in(0,1)$，$\nu_x$ 为分布于 $[-1,1]$ 上的概率测度，其密度 $f(y):=(\pi\sqrt{1-y^2})^{-1}$（$y\in(-1,1)$）。若 $\nu$ 与某个 $v\in L^1(0,1;\mathbb{R})$ 相伴，则由\cref{proposition:ch05-Dirac测度}，对 $\mu$-几乎所有 $x\in(0,1)$ 有 $\nu_x=\delta_{v(x)}$。但 $\nu_x$ 对 $\mu$-几乎所有 $x$ 为非原子的，故矛盾。证明完成。
\end{proof}

\subsection{$\{u_j\}$ 的一致可积性}

我们证明 $\{u_j\}\subseteq L^1(U;\mathbb{R}^d)$ 的一致可积性可由相应 Young 测度序列的窄收敛导出。回忆一致可积性模
\begin{align}
\eta(\{u_j\})(t):=\sup_{j\in\mathbb{N}}\int_{\{|u_j|>t\}}|u_j(x)|\,\mathrm{d}\mu(x)
\qquad(t>0)\label{align:ch05-5.68}
\end{align}
（见\cref{definition:ch02-一致可积}）。

\begin{lemma}\label{lemma:ch05-一致可积}
设 $\{u_j\}\subseteq L^1(U;\mathbb{R}^d)$，$\nu_{u_j}\stackrel{n}{\rightharpoonup}\nu$。则对任意 $t>0$，
\begin{align}
\lim_{j\to\infty}\int_{\{|u_j|>t\}}|u_j(x)|\,\mathrm{d}\mu(x)
=\int_U\mathrm{d}\mu(x)\int_{\{|y|>t\}}|y|\,\mathrm{d}\nu_x(y).\label{align:ch05-5.69}
\end{align}
\end{lemma}

\begin{proof}
设 $\psi_t:\mathbb{R}^d\to\mathbb{R}$，$\psi_t(y):=|y|\chi_{\{|y|>t\}}$。映射 $(x,y)\mapsto\psi_t(y)$ 不是连续的，但 $\psi_t$ 可用连续函数一致逼近，因为由假设 $\{\nu_{u_j}\}$ 窄收敛。更确切地，对任意 $\epsilon>0$，存在 $\varphi_t\in C(\mathbb{R}^d)$ 使
\begin{align*}
\sup_{y\in\mathbb{R}^d}|\psi_t(y)-\varphi_t(y)|<\epsilon.
\end{align*}
于是由\cref{definition:ch05-Young窄}，
\begin{align*}
\lim_{j\to\infty}\int_U\varphi_t(u_j(x))\,\mathrm{d}\mu(x)
=\lim_{j\to\infty}\int_{U\times\mathbb{R}^d}\varphi_t(y)\,\mathrm{d}\nu_{u_j}
=\int_U\mathrm{d}\mu(x)\int_{\mathbb{R}^d}\varphi_t(y)\,\mathrm{d}\nu_x(y),
\end{align*}
且
\begin{align*}
\int_U\varphi_t(u_j(x))\,\mathrm{d}\mu(x)
-\epsilon\mu(U)
\leqslant\int_{\{|u_j|>t\}}|u_j(x)|\,\mathrm{d}\mu(x)
\leqslant\int_U\varphi_t(u_j(x))\,\mathrm{d}\mu(x)
+\epsilon\mu(U).
\end{align*}
令 $j\to\infty$ 与 $\epsilon\to0$ 即得\eqref{align:ch05-5.69}。
\end{proof}

\begin{proposition}\label{proposition:ch05-biting}\label{proposition:ch05-5.4.10}
设 $\{u_j\}\subseteq L^1(U;\mathbb{R}^d)$，$\nu_{u_j}\stackrel{n}{\rightharpoonup}\nu$。假设对每个 $h\in C(\mathbb{R}^d)$ 满足 $h(y)/(1+|y|)\to0$ 当 $|y|\to\infty$，序列 $\{h\circ u_j\}$ 在 $L^1(U)$ 中有界且一致可积。则 $h\circ u_{j_k}\rightharpoonup h_*$ 在 $L^1(U)$ 中，其中子列 $\{u_{j_k}\}\subseteq\{u_j\}$ 使 $\nu_{u_{j_k}}\stackrel{n}{\rightharpoonup}\nu$，且
\begin{align}
h_*(x):=\int_{\mathbb{R}^d}h(y)\,\mathrm{d}\nu_x(y)
\quad\text{对 }\mu\text{-几乎所有 }x\in U.\label{align:ch05-5.70}
\end{align}
\end{proposition}

\begin{proof}
设 $g\in L^\infty(U;\mathbb{R}^d)$。固定任意 $h$ 如上，对任意 $\epsilon>0$，存在 $r>0$ 使
\begin{align*}
\sup_{j\in\mathbb{N}}\int_{\{|u_j|>r\}}|h(u_j(x))|\,\mathrm{d}\mu(x)
=\eta(\{h\circ u_j\})(r)<\epsilon.
\end{align*}
取 $\chi_r\in C_c(\mathbb{R}^d)$ 使 $\chi_r=1$ 在 $\{|y|\leqslant r\}$ 上。则
\begin{align*}
\left|\int_U g[h\circ u_j-h_*]\,\mathrm{d}\mu\right|
\leqslant\left|\int_U g\,[h\chi_r]\circ u_j\,\mathrm{d}\mu-\int_U g h_*\chi_r\,\mathrm{d}\mu\right|
+\left|\int_U g[h(1-\chi_r)]\circ u_j\,\mathrm{d}\mu\right|
+\left|\int_U g h_*(1-\chi_r)\,\mathrm{d}\mu\right|\\
=:I_1+I_2+I_3.
\end{align*}
由\cref{definition:ch05-Young窄}、\eqref{align:ch05-5.66}与\cref{corollary:ch05-Young积分}，
\begin{align*}
I_1=\left|\int_U g(x)\int_{\mathbb{R}^d}h(y)\chi_r(y)\,\mathrm{d}\nu_{u_j}(x,y)
-\int_U g(x)\int_{\mathbb{R}^d}h(y)\chi_r(y)\,\mathrm{d}\nu_x(y)\right|.
\end{align*}
于是由 $\nu_{u_j}\stackrel{n}{\rightharpoonup}\nu$，$I_1\to0$。此外，
\begin{align*}
I_2\leqslant\|g\|_\infty\int_{\{|u_j|>r\}}|h(u_j(x))|\,\mathrm{d}\mu(x)
\leqslant\|g\|_\infty\epsilon,
\end{align*}
且由\cref{lemma:ch05-一致可积}，
\begin{align*}
I_3\leqslant\|g\|_\infty\int_{\{|y|>r\}}|h(y)|\,\mathrm{d}\nu_x(y)
\leqslant\|g\|_\infty\liminf_{j\to\infty}\int_{\{|u_j|>r\}}|h(u_j(x))|\,\mathrm{d}\mu(x)
\leqslant\|g\|_\infty\epsilon.
\end{align*}
综上，
\begin{align*}
\limsup_{j\to\infty}\left|\int_U g[h\circ u_j-h_*]\,\mathrm{d}\mu\right|
\leqslant2\|g\|_\infty\epsilon.
\end{align*}
令 $\epsilon\to0$ 即得结论。
\end{proof}

\begin{remark}\label{remark:ch05-biting注}
\label{remark:ch05-5.4.4}
在\cref{proposition:ch05-biting}的相同假设下，若 $h$ 还满足 $|h(y)|\leqslant C(1+|y|)$（对某个 $C>0$），则 $\{h\circ u_j\}$ 在 $L^1(U)$ 中有界。
\end{remark}

\subsection{Biting 引理}\label{subsection:ch05-Biting}

我们证明如下结果（见\cite[Theorem 2.1]{ref92}）。

\begin{theorem}[Biting 引理]\label{theorem:ch05-biting定理}\label{proposition:ch05-5.4.13}
设 $\{u_j\}\subseteq L^1(U;\mathbb{R}^d)$ 有界。则存在子列 $\{u_{j_k}\}\subseteq\{u_j\}$、递减可测集序列 $\{E_m\}\subseteq\mathscr{B}(U)$（$\mu(\bigcap_{m=1}^\infty E_m)=0$）与 Young 测度 $\nu\in\mathcal{Y}^+(U;\mathbb{R}^d)$ 使对每个固定的 $m\in\mathbb{N}$，序列 $\{u_{j_k}\chi_{E_m^c}\}_{k\in\mathbb{N}}$ 一致可积。
\end{theorem}

\begin{proof}
对任意 $j\in\mathbb{N}$，置 $F_j(t):=\int_{\{|u_j|>t\}}|u_j(x)|\,\mathrm{d}\mu(x)$（$t>0$）。显然对所有 $j$ 与 $t$ 有 $F_j(t)\to0$ 当 $t\to\infty$。由于 $\{u_j\}$ 在 $L^1(U;\mathbb{R}^d)$ 中有界，定义\eqref{align:ch05-5.68}中的一致可积性模 $\eta(\{u_j\})(t)$ 满足 $\eta(\{u_j\})(t)\to0$ 当 $t\to\infty$（见\cite[Theorem 1.29]{ref42}）。于是存在递增序列 $t_m\to\infty$（$m\in\mathbb{N}$）使对所有 $m$，
\begin{align}
\eta(\{u_j\})(t_m)\to0\qquad(m\to\infty).\label{align:ch05-5.71}
\end{align}
对任意固定的 $m$，置 $E_{j,m}:=\{|u_j|>t_m\}$。则由\eqref{align:ch05-5.71}，存在子列 $\{u_{j_k}\}$（由对角论证选取，独立于 $m$）使对每个固定的 $m$，
\begin{align*}
\sup_{k\in\mathbb{N}}\int_{E_{j_k,m}}|u_{j_k}(x)|\,\mathrm{d}\mu(x)\to0.
\end{align*}
由\cref{theorem:ch05-ball}，可进一步选取子列使 $\nu_{u_{j_k}}\stackrel{*}{\rightharpoonup}\nu$。对每个 $m$，置 $E_m:=\limsup_{k\to\infty}E_{j_k,m}$。则 $\{E_m\}$ 递减且 $\mu(\bigcap_m E_m)=0$。对固定 $m$，序列 $\{u_{j_k}\chi_{E_m^c}\}$ 一致可积，因为 $|u_{j_k}|\leqslant t_m$ 在 $E_m^c\setminus E_{j_k,m}$ 上，而在 $E_{j_k,m}^c$ 上 $|u_{j_k}|\leqslant t_m$。证明完成。
\end{proof}

由 Biting 引理与 Young 测度方法，我们得到有界序列 $\{u_j\}\subseteq L^1(U;\mathbb{R}^d)$ 的如下分解。

\begin{theorem}\label{theorem:ch05-分解}\label{proposition:ch05-集中测度极限}
设 $\{u_j\}\subseteq L^1(U;\mathbb{R}^d)$ 有界，$\nu_{u_{j_k}}\stackrel{*}{\rightharpoonup}\nu$。则存在子列 $\{u_{j_k}\}$、递减集序列 $\{E_m\}\subseteq\mathscr{B}(U)$（$\mu(\bigcap_m E_m)=0$）与 $u_*\in L^1(U;\mathbb{R}^d)$ 使
\begin{enumerate}[(i)]
\item $u_{j_k}\chi_{E_m^c}\rightharpoonup u_*\chi_{E_m^c}$ 在 $L^1(U;\mathbb{R}^d)$ 中对每个 $m$；
\item $u_{j_k}\chi_{E_k^c}\rightharpoonup u_*$ 在 $L^1(U;\mathbb{R}^d)$ 中；
\item 存在正 Radon 测度 $\sigma\in\mathscr{R}^+(U;\mathbb{R}^d)$ 使 $u_{j_k}\stackrel{*}{\rightharpoonup}u_*+\sigma$ 在 $(L^\infty(U;\mathbb{R}^d))^*$ 中，且 $\sigma\perp\mu$。
\end{enumerate}
\end{theorem}

\begin{proof}
由\cref{theorem:ch05-biting定理}，存在子列 $\{u_{j_k}\}$、递减集序列 $\{E_m\}$（$\mu(\bigcap_m E_m)=0$）使对每个 $m$，$\{u_{j_k}\chi_{E_m^c}\}$ 一致可积。则由\cref{proposition:ch05-biting}（取 $h(y)=y$），$u_{j_k}\chi_{E_m^c}\rightharpoonup u_*\chi_{E_m^c}$ 在 $L^1(U;\mathbb{R}^d)$ 中，其中 $u_*(x):=\int_{\mathbb{R}^d}y\,\mathrm{d}\nu_x(y)$。这证明了 (i)。

(ii) 由 (i) 与对角论证得到。(iii) 由 (ii) 与测度分解得到：弱*极限 $u_{j_k}\stackrel{*}{\rightharpoonup}u_*+\sigma$ 中的奇异部分 $\sigma$ 集中在 $\bigcap_m E_m$ 上，故 $\sigma\perp\mu$。证明完成。
\end{proof}

作为 Biting 引理的推论，我们还得到如下结果。

\begin{proposition}\label{proposition:ch05-函数弱星}
设 $\{u_j\}\subseteq L^1(U;\mathbb{R}^d)$ 有界，$\nu_{u_{j_k}}\stackrel{*}{\rightharpoonup}\nu$。设 $f:\mathbb{R}^d\to\mathbb{R}$ 连续且 $|f(y)|\leqslant C(1+|y|)$。则存在 $\sigma_f\in\mathscr{R}^+(U)$ 使 $f\circ u_{j_k}\stackrel{*}{\rightharpoonup}f_*+\sigma_f$ 在 $(L^\infty(U))^*$ 中，其中
\begin{align*}
f_*(x):=\int_{\mathbb{R}^d}f(y)\,\mathrm{d}\nu_x(y).
\end{align*}
\end{proposition}

\begin{proof}
由\cref{theorem:ch05-分解}，$u_{j_k}=u_*\chi_{E_k^c}+$ 集中在 $E_k$ 上的部分。对 $f\circ u_{j_k}$ 应用相同分解，其绝对连续部分极限为 $f_*$，奇异部分记为 $\sigma_f$。结论得证。
\end{proof}

最后，我们证明强收敛的如下刻画。

\begin{proposition}\label{proposition:ch05-强收敛}
设 $\{u_j\}\subseteq L^1(U;\mathbb{R}^d)$ 有界，$u_j\to u$ 对 $\mu$-几乎所有 $x\in U$，且 $\|u_j\|_{L^1(U)}\to\|u\|_{L^1(U)}$。则 $u_j\to u$ 在 $L^1(U;\mathbb{R}^d)$ 中强收敛。
\end{proposition}

\begin{proof}
由假设，$u_j(x)\to u(x)$ 对 $\mu$-几乎所有 $x$，故 $|u_j(x)|\to|u(x)|$ 且 $|u_j(x)|-|u(x)|\to0$。由于 $\|u_j\|_{L^1}\to\|u\|_{L^1}$，映射 $x\mapsto||u_j(x)|-|u(x)||$ 一致可积（例如见\cite[Theorem 1.29]{ref42}）。于是由控制收敛定理，
\begin{align*}
\|u_j-u\|_{L^1(U)}=\int_U|u_j(x)-u(x)|\,\mathrm{d}\mu(x)\to0.
\end{align*}
证明完成。
\end{proof}

\end{document}
